\providecommand{\costfiles}[1]{\begin{itemize}[nosep,leftmargin=1.5em,label=--]\footnotesize #1\end{itemize}}
\section{Torsion specialization: selection theorems}\label{sec:tors}

\providecommand{\torsHs}{H_{*}}
\providecommand{\torslen}{\operatorname{length}}
\providecommand{\torscoker}{\operatorname{coker}}
\providecommand{\torsadj}{\operatorname{adj}}
\providecommand{\torsFitt}{\operatorname{Fitt}}
\providecommand{\torsZ}[2]{Z_{[#1,#2]}}
\providecommand{\torsU}{\mathcal{U}}
\providecommand{\torsfiles}[1]{\begin{itemize}[nosep,leftmargin=1.5em,label=--]\footnotesize #1\end{itemize}}

Fix a prime $p$. An integral $p$-adic polynomial or a bounded power series is evaluated at
\emph{torsion points}, that is, at tuples $(\zeta_1-1,\dots,\zeta_s-1)$ in which every
$\zeta_j$ is a root of unity of $p$-power order. The source proof uses such evaluations to
construct its anticyclotomic twists (\cref{src:sec:mechanism}): finitely many integral
polynomials are given, each with a witness point in the open unit polydisc at which its
valuation is known, and one needs a single torsion tuple at which every polynomial has
almost the valuation it has at its witness. The upper bound $H$ for the order exponents of
the $\zeta_j$ has to be fixed before the polynomials are supplied and must not depend on
their degrees. The source obtains $H$ from a divergent harmonic sum.

This section proves the existence theorems of this kind. They are local statements,
independent of \cref{src:thm:dgrh}, and they differ in what is compared and over which
coefficients.
\begin{itemize}
\item For one function, or for the \emph{total} valuation of a family, an exact
  comparison with the Gauss valuation at a fixed positive radius holds with the upper
  exponent $h_0+\floor{Bp^{h_0-1}}$, where $B$ is the budget (\cref{tors:thm:selection}).
  This exponent is optimal by \cref{cost:thm:minimax}.
\item If only nonvanishing is required, the interval of exponents shrinks to
  $[h_*(r),h_*(r)+\floor{B}]$ (\cref{tors:thm:nonvanishing}). More generally, the exact
  budget needed to force a zero at one of a prescribed set of unit powers on a shell is a
  finite covering density, in every dimension (\cref{tors:thm:cover}); for the first two
  powers this gives the optimal endpoint, except that for odd $\ord_p(2)$ and $2B\notin\Z$ it
  is known only up to one (\cref{tors:thm:twopowers}).
\item \emph{Separate} bounds for the members of a family, with unrestricted integral
  coefficients, require a positive tolerance. For a family of fixed cardinality the
  exponent is polynomial in the inverse tolerance, and for polynomials with algebraic
  integral coefficients a positive fraction of the primitive roots of one order is good
  (\cref{tors:thm:fraction,tors:thm:packet}). For families of
  arbitrary cardinality with a finite total budget the harmonic bound of the source
  extends to bounded power series over $\OO_{\Cp}$ (\cref{tors:thm:uniform}).
\item Over a fixed finite extension of $\Qp$ separate bounds hold with tolerance zero, in
  one variable with an exponent that depends only on the wild part of the ramification
  index (\cref{tors:thm:fixedfield}) and in several variables without an iterated
  exponential tower (\cref{tors:thm:notower}).
\end{itemize}
\Cref{tors:tab:overview} lists the bounds in a common notation. The consequences for
simultaneous powers, for determinants and Fitting ideals, and for packets of characters
in the arithmetic twisting construction are derived once, in
\cref{tors:sec:powers,tors:sec:matrices,tors:sec:characters}. Lower bounds, exact costs and
the dependence on the coefficient field are the subject of \cref{sec:cost}.

\subsection{Notation, Gauss valuations and the objectives}\label{tors:sec:setup}
\sources{grh_quantitative_twists_2026-11-25; grh_method_transfers_2026-11-25_1; grh_beyond_dirichlet_grh_2026-11-25; grh_exact_torsion_transfer_2026-11-25; grh_fixed_field_transfer_2026-11-25; grh_sharp_transfers_2026-11-25}

Let $v$ be the valuation of $\Cp$ with $v(p)=1$ and $v(0)=+\infty$, and let
$\OO=\OO_{\Cp}=\set{x\in\Cp: v(x)\ge0}$. Since $\Cp$ is algebraically closed, every
nonzero polynomial over $\Cp$ is a constant times a product of linear factors; the
root-counting arguments below therefore need no algebraicity assumption on the
coefficients. Write $\overline{\Q}_p$ for the algebraic closure of $\Qp$ in $\Cp$ and
$\overline{\Z}_p=\overline{\Q}_p\cap\OO$ for its ring of integral elements. Some selection
theorems below are proved, as in their sources, for polynomials with coefficients in
$\overline{\Z}_p$ and witnesses with algebraic coordinates; their extensions to series
over $\OO$ go through truncation and algebraic approximation
(\cref{tors:lem:truncation,tors:lem:approx}).

\paragraph{Shells.}
For $h\ge1$ put
\[
  e_h=\varphi(p^h)=(p-1)p^{h-1},\qquad r_h=\frac1{e_h},\qquad
  Z_h=\set{\zeta-1:\ \zeta\text{ a root of unity of exact order }p^h}.
\]
The set $Z_h$, the $h$-th \emph{shell}, has $e_h$ elements, and
\begin{equation}\label{tors:eq:shell}
  v(z)=r_h\qquad(z\in Z_h),
\end{equation}
because $\Phi_{p^h}(1)=p$ is the product of the $e_h$ conjugate numbers $1-\zeta$, which
have equal valuations. If $\xi$ is a primitive $p^h$-th root of unity and $a\not\equiv b
\pmod{p^h}$ are prime to $p$, then $\xi^{a-b}$ has order $p^{h-\vp(a-b)}$, so
\begin{equation}\label{tors:eq:distance}
  v(\xi^a-\xi^b)=r_{h-\vp(a-b)}=p^{\vp(a-b)}\,r_h .
\end{equation}
For a real $\rho>0$ let
\[
  h_*(\rho)=\min\set{h\ge1:\ r_h\le\rho}.
\]
For integers $1\le h_0\le H$ put $\torsZ{h_0}{H}=\bigcup_{h_0\le h\le H}Z_h$. Thus
$z\in\torsZ{h_0}{H}$ if and only if $z=\zeta-1$ with $\zeta^{p^H}=1$ and
$r_H\le v(\zeta-1)\le r_{h_0}$, and $\#\torsZ{h_0}{H}=p^H-p^{h_0-1}$. The set
$\torsZ{h_0}{H}^s$ is called the \emph{torsion grid}; its points are the torsion tuples
whose coordinates have order exponents in $[h_0,H]$. The coordinates of one tuple may have
different orders.

\paragraph{Gauss valuations.}
A \emph{bounded series} is a formal power series $F=\sum_\nu a_\nu T^\nu\in
\Cp[[T_1,\dots,T_s]]$ whose coefficient valuations are bounded below; it is \emph{integral}
if all $a_\nu\in\OO$. For a vector of weights $\rho=(\rho_1,\dots,\rho_s)$ with all
$\rho_j>0$, and a nonzero bounded series $F$, put
\begin{gather*}
  G_\rho(F)=\min_\nu\bigl\{v(a_\nu)+\rho_1\nu_1+\dots+\rho_s\nu_s\bigr\},\\
  G_0(F)=\inf_\nu v(a_\nu),\qquad \Delta_\rho(F)=G_\rho(F)-G_0(F)\ge0 .
\end{gather*}
The minimum exists because the $v(a_\nu)$ are bounded below and the weights are positive;
the infimum $G_0(F)$ need not be attained. When all weights are equal to one number
$\rho>0$ we write $G_\rho$ for that number, so $G_\rho(F)=\min_\nu\{v(a_\nu)+\rho\abs{\nu}\}$
with $\abs\nu=\nu_1+\dots+\nu_s$. For a polynomial in one variable the definition of
$G_u$ also makes sense at $u=0$ and agrees with $G_0$. The quantity $\Delta_\rho(F)$ is
unchanged when $F$ is multiplied by a nonzero scalar, and $\Delta_\rho(F)\le G_\rho(F)$ when
$F$ is integral. Lowering a weight does not increase $G_\rho(F)$.

\begin{lemma}[Gauss valuations]\label{tors:lem:gauss}
Let $F,G$ be nonzero bounded series in the same variables and let $\rho$ be a vector of
positive weights.
\begin{enumerate}[label=\textup{(\alph*)}]
\item $G_\rho(F+G)\ge\min\{G_\rho(F),G_\rho(G)\}$ and $G_\rho(FG)=G_\rho(F)+G_\rho(G)$.
  Consequently $G_0$ and $\Delta_\rho$ are additive on products.
\item If $x\in\Cp^s$ and $v(x_j)\ge\rho_j$ for all $j$, then $F(x)$ converges and
  $v(F(x))\ge G_\rho(F)$.
\item If $g=b\prod_\alpha(T-\alpha)\in\Cp[T]$, with the roots counted with multiplicity,
  then for every real $u\ge0$
  \begin{equation}\label{tors:eq:gauss}
    G_u(g)=v(b)+\sum_\alpha\min\{u,v(\alpha)\}.
  \end{equation}
  The formula includes roots of negative valuation and the root $0$.
\end{enumerate}
\end{lemma}

\begin{proof}
(a) The first inequality is clear, and every term of the product $FG$ has weighted
valuation at least $G_\rho(F)+G_\rho(G)$. Among the exponents attaining the minimum for
$F$ let $\nu_0$ be the lexicographically smallest, and define $\mu_0$ in the same way for
$G$. The coefficient of $T^{\nu_0+\mu_0}$ in $FG$ is a finite sum of products
$a_\nu b_\mu$ with $\nu+\mu=\nu_0+\mu_0$. If $(\nu,\mu)\ne(\nu_0,\mu_0)$, then $\nu$ is
lexicographically smaller than $\nu_0$ or $\mu$ is lexicographically smaller than $\mu_0$,
so one of the two factors does not attain its minimum and the weighted valuation of
$a_\nu b_\mu$ is strictly larger than $G_\rho(F)+G_\rho(G)$. Hence the term
$a_{\nu_0}b_{\mu_0}$ cannot cancel, and $G_\rho(FG)=G_\rho(F)+G_\rho(G)$. As all weights
tend to $0$ the Gauss valuations decrease to $G_0$, which is therefore additive as well.

(b) Each term satisfies $v(a_\nu x^\nu)\ge v(a_\nu)+\sum_j\rho_j\nu_j\ge G_\rho(F)$, and
$v(a_\nu x^\nu)\ge G_0(F)+\abs\nu\min_j\rho_j\to\infty$. Completeness of $\Cp$ and the
ultrametric inequality give the claim.

(c) This follows from (a), which for polynomials in one variable holds also at $u=0$ by
the same argument, since $G_u(T-\alpha)=\min\{u,v(\alpha)\}$.
\end{proof}

A \emph{witness of depth $\rho$} for an integral series $F$ is a point $x\in\Cp^s$ with
$v(x_j)\ge\rho$ for all $j$ and $v(F(x))<\infty$. By \cref{tors:lem:gauss}(b),
\begin{equation}\label{tors:eq:witness}
  G_\rho(F)\le v(F(x)) .
\end{equation}
A budget for witness valuations is therefore also a budget for the Gauss valuations at
radius $\rho$, and a hypothesis stated with $G_\rho$ is the weaker one. The number
$v(F(x))$ is also called an interior witness valuation, or a reference value.

\begin{lemma}[Roots on the shells]\label{tors:lem:shells}
Let $g=b\prod_\alpha(T-\alpha)\in\Cp[T]$ be nonzero and let $m_h(g)$ be the number of
roots of valuation exactly $r_h$, counted with multiplicity. Let $\rho>0$.
\begin{enumerate}[label=\textup{(\alph*)}]
\item $\displaystyle\sum_{h:\ r_h\le\rho}r_h\,m_h(g)\le\Delta_\rho(g)$, and
  $\Delta_\rho(g)\le G_\rho(g)$ if $g$ is integral.
\item For every $h\ge1$ and $z\in Z_h$,
  \begin{equation}\label{tors:eq:value}
    v(g(z))=G_{r_h}(g)+\sum_{v(\alpha)=r_h}\bigl(v(z-\alpha)-r_h\bigr),
  \end{equation}
  and every summand is nonnegative or $+\infty$.
\item If $r_h\le\rho$, then
  \begin{equation}\label{tors:eq:slack}
    G_\rho(g)-G_{r_h}(g)\ \ge\sum_{k<h,\ r_k\le\rho}(r_k-r_h)\,m_k(g)\ \ge0 .
  \end{equation}
\end{enumerate}
\end{lemma}

\begin{proof}
By \eqref{tors:eq:gauss}, $\Delta_\rho(g)=\sum_\alpha\bigl(\min\{\rho,v(\alpha)\}-
\min\{0,v(\alpha)\}\bigr)$. Every summand is nonnegative, and a root of valuation
$r_h\le\rho$ contributes exactly $r_h$; distinct $h$ have distinct radii. This proves (a).
For (b), $v(g(z))=v(b)+\sum_\alpha v(z-\alpha)$, and $v(z-\alpha)=\min\{r_h,v(\alpha)\}$
when $v(\alpha)\ne r_h=v(z)$, while $v(z-\alpha)\ge r_h$ otherwise; compare with
\eqref{tors:eq:gauss} at $u=r_h$. For (c), \eqref{tors:eq:gauss} gives
$G_\rho(g)-G_{r_h}(g)=\sum_\alpha\bigl(\min\{\rho,v(\alpha)\}-\min\{r_h,v(\alpha)\}\bigr)$
with nonnegative summands, and a root of valuation $r_k$ with $r_h<r_k\le\rho$
contributes $r_k-r_h$.
\end{proof}

Inequality \eqref{tors:eq:slack} is the \emph{saving from earlier shells}: a root on an
earlier shell has not only used a part of the budget, it also lowers the target
$G_{r_h}(g)$ on every later shell relative to the fixed target $G_\rho(g)$.

\paragraph{The objectives.}
Let $F_i$ be nonzero integral series with real \emph{targets} $t_i\ge G_\rho(F_i)$, for
instance $t_i=v(F_i(x_i))$ for witnesses $x_i$ of depth $\rho$; the number
$\sum_it_i$ is the \emph{budget}. At a torsion tuple $z$ the following comparisons occur.
\begin{itemize}
\item \emph{Separate comparison with tolerances $\eps_i\ge0$:} $v(F_i(z))\le t_i+\eps_i$
  for every $i$. A stronger form bounds the \emph{total positive excess}
  $\sum_i\bigl(v(F_i(z))-t_i\bigr)_+$, where $(u)_+=\max\{u,0\}$.
\item \emph{Aggregate comparison:} $\sum_iv(F_i(z))\le\sum_iG_\rho(F_i)$. By
  \cref{tors:lem:gauss}(a) this is the comparison $v(F(z))\le G_\rho(F)$ for the single
  series $F=\prod_iF_i$.
\item \emph{Nonvanishing:} $F_i(z)\ne0$ for every $i$.
\end{itemize}
In the first two comparisons the radius $\rho$ is fixed. The harmonic lemma of the source
compares $v(g(z))$, for $z\in Z_h$, with the \emph{moving} target $G_{r_h}(g)\le G_\rho(g)$,
which is a stronger requirement. The aggregate comparison does not imply the separate
comparisons with tolerance zero: a negative deviation of one factor from its Gauss
valuation can compensate a positive deviation of another. With unrestricted coefficients
in $\OO$ and a bounded budget, neither separate exact comparisons for families of unbounded
cardinality nor an exact comparison with the moving target can be guaranteed on any
prescribed finite set of shells; the obstructions are constructed in \cref{sec:cost}.
This is why tolerance zero appears below only for the aggregate comparison
(\cref{tors:sec:optimal}) and over a fixed coefficient field (\cref{tors:sec:fixedfield}),
and why the separate comparisons over $\OO$ carry a positive tolerance
(\cref{tors:sec:tolerance}).

\begin{table}[ht]
\centering\small
\renewcommand{\arraystretch}{1.25}
\begin{tabular}{@{}p{0.16\linewidth}p{0.22\linewidth}p{0.26\linewidth}p{0.27\linewidth}@{}}
\toprule
Result & Data & Conclusion at one torsion tuple $z$ & Sufficient upper exponent $H$\\
\midrule
\Cref{tors:thm:selection}
 & bounded series over $\Cp$, any $s$, $\Delta_\rho(F)\le B$; or a finite family with
   $\sum_i\Delta_\rho(F_i)\le B$
 & $v(F(z))\le G_\rho(F)$; for a family
   $\sum_iv(F_i(z))\le\sum_iG_\rho(F_i)$
 & $h_0+\floor{Bp^{h_0-1}}$ (optimal)\\
\Cref{tors:thm:nonvanishing}
 & finite family over $\OO$, any $s$, witnesses of depth $r$, budget $B$
 & $F_i(z)\ne0$ for all $i$
 & $h_*(r)+\floor{B}$ (optimal)\\
\Cref{tors:thm:twopowers}
 & as above, $p$ odd, $r\ge r_{h_0}$
 & $F_i(z)\ne0$ and $F_i([2]z)\ne0$ for all $i$
 & $h_0+\floor{2B}$, or $h_0+\floor{2B}-1$ if $\ord_p(2)$ is odd and $2B\in\Z_{>0}$ (optimal,
   except for odd $\ord_p(2)$ and $2B\notin\Z$, where it is exact up to one)\\
\Cref{tors:thm:fraction}
 & $n$ polynomials over $\overline{\Z}_p$, $G_\rho(g_i)\le B_i$, tolerances $\delta_i>0$
 & $v(g_i(z))\le G_\rho(g_i)+\delta_i$ for all $i$, for more than a fraction $1-\eta$ of
   one shell
 & $\ceil[\big]{(h_0+C)P^{1/(c\eta)}}$, with $C,P$ as in \eqref{tors:eq:CP}\\
\Cref{tors:prop:linear}
 & one polynomial over $\OO$, $G_\rho(g)\le B$, tolerance $\delta>0$
 & $v(g(z))\le G_\rho(g)+\delta$
 & $h_0+\ceil{Bt_0/(\kappa\delta)}$\\
\Cref{tors:thm:uniform}
 & any family over $\OO$, $s\le s_0$ variables, witnesses of depth $\rho$, budget $B$
 & $\sum_i\bigl(v(F_i(z))-t_i\bigr)_+\le\eps$
 & $\ceil[\big]{(h_*(\rho)+1)\exp\bigl((B+1)(s_0+2)/\min\{\eps,1\}\bigr)}$\\
\Cref{tors:thm:fixedfield}
 & finite family over $\OO_K$, one variable, $\sum_it_i\le B$
 & $v(g_i(z))\le t_i$ for all $i$
 & $h_0+\floor{p^{w}B}$, $w=\vp(e(K/\Qp))$\\
\Cref{tors:thm:notower}
 & finite family over $\OO_K$, $d$ variables, $\sum_iG_\rho(F_i)\le B$ (series and
   infinite families: \cref{tors:cor:notowerseries})
 & $v(F_i(z))\le G_\rho(F_i)$ for all $i$
 & $h_0+\ceil[\big]{(3^{d-1}-1)\log_p(1+B/\rho)}+\floor{p^wB}$\\
\bottomrule
\end{tabular}
\caption{The selection theorems in one notation. In every row $r_{h_0}\le\rho$, the
coordinates of $z$ have order exponents in $[h_0,H]$ (with $h_0=h_*(\rho)$ or $h_*(r)$
where no $h_0$ is named), and $H$ depends neither on degrees nor, except in the last two
rows, on the coefficient field.}
\label{tors:tab:overview}
\end{table}

The passage from polynomials to power series uses the finiteness of the torsion grid and
the fact that $H$ is known in advance.

\begin{lemma}[Truncation]\label{tors:lem:truncation}
Let $F$ be a nonzero bounded series in $s$ variables, let $a\le G_0(F)$ and $M$ be real
numbers, and let $r>0$. Choose an integer $D\ge0$ with $a+(D+1)r>M$, and let $P$ be the
sum of the terms of $F$ of total degree at most $D$.
\begin{enumerate}[label=\textup{(\alph*)}]
\item If $y\in\Cp^s$ has all coordinates of valuation at least $r$, then
  $v(F(y)-P(y))>M$. In particular $v(P(y))\le M$ implies $v(F(y))=v(P(y))$, and
  $v(F(y))\le M$ implies $v(P(y))=v(F(y))$.
\item If $\rho\ge r$ and $G_\rho(F)\le M$, then $P\ne0$ and $G_\rho(P)=G_\rho(F)$.
  Moreover $G_0(P)\ge a$.
\end{enumerate}
\end{lemma}

\begin{proof}
Every omitted term $a_\nu T^\nu$ has $\abs\nu\ge D+1$ and $v(a_\nu)\ge a$. At $y$ its
valuation is at least $a+(D+1)r>M$, and its weighted valuation at radius $\rho\ge r$ is
larger than $M\ge G_\rho(F)$. The ultrametric inequality gives both parts.
\end{proof}

\begin{lemma}[Algebraic approximation]\label{tors:lem:approx}
Let $P$ be a polynomial in $s$ variables with coefficients in $\OO$, let $x\in\OO^s$ and
let $M>0$. There are a polynomial $Q$ with coefficients in $\overline{\Z}_p$ and a point
$x'\in\overline{\Z}_p^{\,s}$ such that $v(P(y)-Q(y))>M$ for every $y\in\OO^s$,
$v(x'_j-x_j)>M$ for every $j$, and $v(Q(x')-P(x))>M$. If $v(x_j)\ge\rho$ with
$\rho\le M$, then $v(x'_j)\ge\rho$.
\end{lemma}

\begin{proof}
Since $\overline{\Q}_p$ is dense in $\Cp$, each of the finitely many coefficients of $P$ and
each coordinate of $x$ can be approximated by an element of $\overline{\Z}_p$ with error of
valuation greater than $M$; a zero coordinate is kept. Replacing the coefficients changes
the value at every point of $\OO^s$ by an amount of valuation greater than $M$. Integral
polynomials are nonexpanding on $\OO^s$: telescoping the difference of two monomials with
integral arguments shows that changing every coordinate by an amount of valuation greater
than $M$ changes the value by an amount of valuation greater than $M$. The last assertion
is the ultrametric inequality.
\end{proof}

\subsection{Setting, root accounting and the two loss conventions}\label{cost:sec:setting}
{\raggedright\sources{grh_sharp_transfers_2026-11-25; grh-torsion-and-prime-sampling; grh-torsion-and-zero-data; grh_fixed_field_transfer_2026-11-25; grh_sharp_torsion_and_symplectic_blocks}}

Fix a prime $p$ and let $v$ be the valuation of $\Cp$ normalized by $v(p)=1$, with
$v(0)=+\infty$. Write $\OO=\OO_{\Cp}$. For $h\ge1$ put
\[
 e_h=(p-1)p^{h-1},\qquad r_h=e_h^{-1},\qquad
 Z_h=\set{\zeta-1:\ord(\zeta)=p^h}.
\]
The \emph{shell} $Z_h$ has $e_h$ points. They are the roots of the Eisenstein polynomial
$\Phi_{p^h}(1+T)$, whose constant term is $p$, so each of them has valuation $r_h$.

For a nonzero polynomial $g(T)=\sum_j a_jT^j\in\Cp[T]$ and a real $u\ge0$ the \emph{Gauss
valuation} and its variation are
\[
 G_u(g)=\min_j\set{v(a_j)+ju},\qquad \Delta_\rho(g)=G_\rho(g)-G_0(g)\quad(\rho\ge0).
\]
If $v(x)\ge u$ then $v(g(x))\ge G_u(g)$. The variation $\Delta_\rho(g)$ is nonnegative and is
unchanged when $g$ is multiplied by a nonzero scalar. If $g$ is integral, then $G_0(g)\ge0$
and $\Delta_\rho(g)\le G_\rho(g)$. Dividing a nonzero polynomial by a coefficient of least
valuation makes it integral with $G_0=0$ and $G_\rho=\Delta_\rho$.

\begin{lemma}[Root accounting]\label{cost:lem:roots}
Let $g=b\prod_\alpha(T-\alpha)$ be the factorization of a nonzero polynomial over $\Cp$, with
multiplicities.
\begin{enumerate}[label=\textup{(\roman*)}]
\item For every $u\ge0$,
\begin{equation}\label{cost:eq:gauss}
 G_u(g)=v(b)+\sum_\alpha\min\set{u,v(\alpha)}.
\end{equation}
This includes roots of negative valuation and the root $0$, whose valuation is $+\infty$.
\item If $v(z)=r$, then
\begin{equation}\label{cost:eq:shellvalue}
 v(g(z))-G_r(g)=\sum_{v(\alpha)=r}\bigl(v(z-\alpha)-r\bigr),
\end{equation}
and every summand is nonnegative.
\item Let $S$ be a finite set of positive integers with $\rho\ge\max_{h\in S}r_h$, and let
$M_h(g)$ be the number of roots of valuation exactly $r_h$. Then
\begin{equation}\label{cost:eq:budget}
 \sum_{h\in S}r_hM_h(g)\le\Delta_\rho(g).
\end{equation}
\end{enumerate}
\end{lemma}

\begin{proof}
The Gauss valuation is multiplicative: in a product, choose in each factor the smallest
exponent attaining the weighted coefficient minimum; the sum of the two exponents is the
smallest exponent of minimum weight in the product, and its minimum-weight contribution
cannot cancel. Since $G_u(T-\alpha)=\min\set{u,v(\alpha)}$, this gives \eqref{cost:eq:gauss}.
For (ii), a root with $v(\alpha)\ne r$ has $v(z-\alpha)=\min\set{r,v(\alpha)}$, which is its
term in \eqref{cost:eq:gauss} at $u=r$; a root with $v(\alpha)=r$ has $v(z-\alpha)\ge r$.
For (iii), by \eqref{cost:eq:gauss},
\[
 \Delta_\rho(g)=\sum_\alpha\bigl(\min\set{\rho,v(\alpha)}-\min\set{0,v(\alpha)}\bigr).
\]
Every summand is nonnegative, a root of valuation $r_h\le\rho$ contributes exactly $r_h$,
and distinct $h$ have distinct radii.
\end{proof}

Two ways of measuring the failure of a specialization occur in the sequel.

\begin{definition}[Aggregate convention]\label{cost:def:aggregate}
Fix $\rho>0$ and $\delta\ge0$. A point $z$ is \emph{good} for a nonzero polynomial $g$ if
$v(g(z))\le G_\rho(g)+\delta$. The polynomial $g$ \emph{blocks} a set $\Sigma$ of points at
tolerance $\delta$ if no point of $\Sigma$ is good, that is,
\[
 v(g(z))>G_\rho(g)+\delta\qquad\text{for every }z\in\Sigma.
\]
The comparison radius $\rho$ is fixed. It is not replaced by the valuation of the point
currently being tested. The \emph{cost} of $g$ is $\Delta_\rho(g)$, which equals $G_\rho(g)$
after the normalization $G_0(g)=0$.
\end{definition}

\begin{definition}[Packets with individual thresholds]\label{cost:def:packet}
Let $R$ be $\OO$ or the ring of integers $\OO_K$ of a finite extension $K/\Qp$ inside $\Cp$,
and fix $\rho>0$. A \emph{packet} over $R$ consists of finitely many nonzero polynomials
$g_i\in R[T]$ and finite thresholds $t_i\ge G_\rho(g_i)$. Its \emph{cost} is $\sum_it_i$,
and its \emph{loss} at a point $z$ is
\[
 \mathcal L(z)=\sum_i\max\set{0,\,v(g_i(z))-t_i}.
\]
Let $S$ be a nonempty finite set of positive integers with $\rho\ge\max_{h\in S}r_h$. The
packet \emph{blocks} $S$ at tolerance $\delta\ge0$ if $\mathcal L(z)>\delta$ for every
$z\in\bigcup_{h\in S}Z_h$. The least cost of a packet over $R$ blocking $S$ at tolerance
$\delta$ is denoted $C_R(S,\delta)$; we write $C_K$ for $C_{\OO_K}$.
\end{definition}

The basic example of a threshold is a witness value: if $v(x_i)\ge\rho$ and
$t_i=v(g_i(x_i))<\infty$, then $t_i\ge G_\rho(g_i)$. All extremal packets below have every
witness equal to $0$ and $t_i=G_\rho(g_i)$, so the minimum costs are the same for witness
thresholds, for Gauss thresholds and for general thresholds. A packet with one member and
$t=G_\rho(g)$ is the aggregate convention for an integral polynomial. For several members the
positive part matters: a saving on one polynomial cannot offset a loss on another. The
aggregate convention applied to the product $\prod_ig_i$ compares
$\sum_iv(g_i(z))$ with $\sum_iG_\rho(g_i)$ and does allow such compensation. The two
conventions have different exact costs (Theorems~\ref{cost:thm:minimax} and~\ref{cost:thm:packet}).

\subsection{Blocking an arbitrary finite set of points}\label{cost:sec:finite}
{\raggedright\sources{grh_method_transfers_research; grh_sharp_torsion_and_symplectic_blocks}}

In the aggregate convention the least cost and the least degree of a blocker can be computed
for every finite set of test points in the open unit disc, at every tolerance. The
computation uses only the valuations of the points and of their differences. All statements
of this subsection are unchanged if $v$, $\rho$ and $\delta$ are multiplied by a common
positive constant.

\begin{lemma}[Greedy sequences]\label{cost:lem:greedy}
Let $X\subset\Cp$ be a set of $l<\infty$ points and $\omega$ a real number with
$v(x-y)\ge\omega$ for all $x,y\in X$. Choose $z_0\in X$ arbitrarily and, for $1\le m<l$,
choose $z_m$ among the unused points $x$ so as to minimize
$\sum_{i<m}(v(x-z_i)-\omega)$. Let $E_X(m)$ be the value of this minimum for $0\le m<l$, so
that $E_X(0)=0$, and put $E_X(l)=+\infty$. Then for $0\le m<l$ every monic polynomial
$Q\in\Cp[T]$ of degree $m$ satisfies
\[
 \min_{x\in X}v(Q(x))\le m\omega+E_X(m),
\]
with equality for $P_m=\prod_{i<m}(T-z_i)$. In particular $E_X(m)$ does not depend on the
choices made, and it is nondecreasing in $m$.
\end{lemma}

\begin{proof}
Put $b_m=m\omega+E_X(m)$. The polynomial $P_m$ vanishes at the used points, so the greedy
choice says $\min_Xv(P_m)=v(P_m(z_m))=b_m$. Write a competing monic polynomial as
$Q=\sum_{i=0}^mc_iP_i$ with $c_m=1$, and suppose $v(Q(x))>b_m$ for all $x\in X$. Evaluate
successively at $z_0,\ldots,z_{m-1}$. At $z_i$ the terms with index larger than $i$ vanish,
and by induction the terms with smaller index have valuation greater than $b_m$ everywhere
on $X$; hence $v(c_i)+b_i>b_m$, and then $v(c_iP_i(x))>b_m$ for every $x\in X$. At $z_m$ all
terms with $i<m$ exceed $b_m$, while $v(P_m(z_m))=b_m$. Thus $v(Q(z_m))=b_m$, a
contradiction. Monotonicity holds because each added summand is nonnegative and the set of
unused points shrinks.
\end{proof}

This is the finite-set argument behind Bhargava's $P$-orderings and generalized
factorials~\cite{Bhargava1997Porderings,Bhargava1998Generalized}.

Now let $\Sigma$ be a nonempty finite subset of $\set{z\in\Cp:v(z)>0}$, and fix $\rho>0$ and
$\delta\ge0$. Give each point the weight $\omega(z)=\min\set{\rho,v(z)}>0$, let
$\Sigma_\omega$ be the set of points of weight $\omega$, and write
$E_\omega=E_{\Sigma_\omega}$ for the function of Lemma~\ref{cost:lem:greedy}. Process the
distinct weights in decreasing order and define integers $m_\omega$ by
\begin{equation}\label{cost:eq:rec}
 t_\omega=\delta+\sum_{u>\omega}m_u(u-\omega),\qquad
 m_\omega=\min\set{0\le m\le\abs{\Sigma_\omega}:E_\omega(m)>t_\omega}.
\end{equation}

\begin{theorem}[Least degree and least budget of a blocker; \stU]\label{cost:thm:finite}
Put $d_*=\sum_\omega m_\omega$ and $B_*=\sum_\omega\omega\,m_\omega$.
\begin{enumerate}[label=\textup{(\alph*)}]
\item Every nonzero $g\in\OO[T]$ that blocks $\Sigma$ at tolerance $\delta$ has, for each
weight $\omega$, at least $m_\omega$ roots $\alpha$ with $\min\set{\rho,v(\alpha)}=\omega$,
counted with multiplicity. In particular $\deg g\ge d_*$ and
$G_\rho(g)\ge\Delta_\rho(g)\ge B_*$.
\item The monic polynomial whose roots are the first $m_\omega$ greedy points of each part
$\Sigma_\omega$ blocks $\Sigma$ at tolerance $\delta$, has degree $d_*$ and has
$G_\rho=B_*$.
\item Consequently every nonzero $g\in\OO[T]$ with $\deg g\le d$ and $G_\rho(g)\le B$ has a
good point in $\Sigma$ if and only if $d<d_*$ or $B<B_*$.
\item Every nonzero $F\in\OO[[T]]$ with $G_\rho(F)<B_*$ has a good point in $\Sigma$.
\end{enumerate}
\end{theorem}

\begin{proof}
Let $g=b\prod_\alpha(T-\alpha)$ be a blocker. By Lemma~\ref{cost:lem:roots}(i), the contribution
of a root $\alpha$ to $v(g(z))-G_\rho(g)$ is $v(z-\alpha)-\min\set{\rho,v(\alpha)}$. Roots
with $v(\alpha)\le0$ contribute zero. For any other root put
$u=\min\set{\rho,v(\alpha)}$. If there is no part of weight $u$, the root contributes at
most zero at every test point and can be discarded. If the part exists, the contribution of
the root at a point of another part of weight $\omega$ is $\min\set{u,\omega}-u$. Within
$\Sigma_u$ the roots of weight $u$ contribute the valuation of their monic product minus
$u$ for each root. This reasoning includes the part of weight $\rho$, which can contain
points of several valuations.

Let $M_u$ be the number of retained roots of weight $u$. At each point of $\Sigma_\omega$,
blocking requires the monic product of the $M_\omega$ roots of weight $\omega$ to have
valuation exceeding $M_\omega\omega$ by strictly more than
$\delta+\sum_{u>\omega}M_u(u-\omega)$; the discarded terms could only have lowered the
excess. By Lemma~\ref{cost:lem:greedy} and induction from the largest weight down,
$M_\omega\ge m_\omega$ for every part; the claim is automatic when
$M_\omega\ge\abs{\Sigma_\omega}$. Each retained root costs its weight in $\Delta_\rho(g)$,
as in the proof of Lemma~\ref{cost:lem:roots}(iii). This proves (a). For the polynomial in (b)
the cross-part effects are exactly the charges in \eqref{cost:eq:rec}, the one-part minima
are attained by Lemma~\ref{cost:lem:greedy}, and $G_0=0$. Hence it blocks, and its Gauss
valuation is $B_*$. Part (c) combines (a) and (b).

For (d), let $A=G_\rho(F)$ and truncate $F$ at a degree $D$ with
$(D+1)\min_{z\in\Sigma}\omega(z)>A+\delta$. The discarded tail has valuation above
$A+\delta$ at each point of $\Sigma$ and leaves the Gauss valuation unchanged, so a good
point of the truncation is a good point of $F$.
\end{proof}

Given the distances, updating one running sum per unused point evaluates the greedy
sequences in $O(\sum_\omega\abs{\Sigma_\omega}^2)$ valuation additions and comparisons. When
the test points are algebraic, the extremal polynomial may require a finite extension of the
base field. Prescribing the field of its coefficients adds Galois constraints and can change
the answer (\cref{cost:sec:fields}).

\subsection{Torsion shells and their capacity}\label{cost:sec:shell}
{\raggedright\sources{grh-torsion-and-zero-data; grh-torsion-and-prime-sampling; grh_sharp_torsion_and_symplectic_blocks; grh_sharp_transfers_2026-11-25}}

The shells $Z_h$ are the first instance of the following setting.

\begin{definition}[Lubin--Tate shells]\label{cost:def:lt}
Let $K/\Qp$ be finite with uniformizer $\pi$ and residue field of cardinality $q$, and let
$F$ be a one-dimensional Lubin--Tate formal $\OO_K$-module for $\pi$, with multiplication
series $[a]$, $a\in\OO_K$, and formal difference $\ominus$. Let $v_F$ be the valuation of
$\Cp$ with $v_F(\pi)=1$, and put
\[
 e_h=(q-1)q^{h-1},\qquad r_h=e_h^{-1},\qquad
 Z_h=\set{x:[\pi^h]x=0,\ [\pi^{h-1}]x\ne0}\quad(h\ge1),
\]
where $x$ ranges over the maximal ideal of $\OO$. The formal parameter itself is evaluated;
there is no subtraction of $1$. We refer to this as setting \textup{(LT)}. Setting
\textup{(C)}, the cyclotomic case of \cref{cost:sec:setting}, is the special case $K=\Qp$,
$\pi=p$, $F=\widehat{\mathbb G}_m$, in which $[a](T)=(1+T)^a-1$, $q=p$ and $v_F=v$.
\end{definition}

\begin{lemma}[The torsion metric]\label{cost:lem:metric}
In setting \textup{(LT)}:
\begin{enumerate}[label=\textup{(\roman*)}]
\item $F[\pi^h]\simeq\OO_K/\pi^h$ as $\OO_K$-modules; the shell $Z_h$ has $e_h$ points, and
each of them has $v_F$-valuation $r_h$.
\item $v_F(x-y)=v_F(x\ominus y)$ for all $x,y$ of positive valuation.
\item Let $t_h$ generate $F[\pi^h]$. Then $Z_h=\set{[a]t_h:a\in(\OO_K/\pi^h)^\times}$, and for
distinct unit residues $a,b$ with $\tau=v_\pi(a-b)\in\set{0,\ldots,h-1}$,
\begin{equation}\label{cost:eq:distance}
 v_F([a]t_h-[b]t_h)-r_h=r_{h-\tau}-r_h=\frac{q^\tau-1}{e_h}=\frac1{e_h}\sum_{j=1}^\tau e_j.
\end{equation}
\end{enumerate}
\end{lemma}

\begin{proof}
The structure of $F[\pi^h]$ is the Lubin--Tate theorem~\cite{LubinTate1965Formal}. The series
$[\pi]$ satisfies $[\pi](T)=\pi T+\cdots$ and $[\pi](T)\equiv T^q\pmod\pi$. Its Newton polygon
shows that a nonzero point killed by $\pi$ has valuation $1/(q-1)$. If $[\pi]x\in Z_{h-1}$
with $h\ge2$, then $v_F(x)<1/(q-1)$, the term $x^q$ is the unique term of least valuation,
and $v_F([\pi]x)=q\,v_F(x)$. Induction gives (i). For (ii), write
$F(X,[-1]Y)=(X-Y)\,U(X,Y)$. The series on the left vanishes at $X=Y$ and has linear term
$X-Y$, so $U$ has integral coefficients and constant term $1$, and $v_F(U)=0$ at points of
positive valuation. For (iii), $[a]t_h\ominus[b]t_h=[a-b]t_h$ has exact order
$\pi^{h-\tau}$; now apply (ii) and (i).
\end{proof}

Thus the unit residues modulo $\pi^h$ form the leaves of a tree with $q-1$ first branches,
in which every later branch has $q$ children, and \eqref{cost:eq:distance} expresses the
distance of two leaves through the depth at which they separate. For $0\le m<e_h$ define the
\emph{capacity}
\begin{equation}\label{cost:eq:capacity}
 U_h(m)=\frac1{e_h}\sum_{j=1}^{h-1}e_j\floor*{\frac m{e_j}},\qquad U_h(e_h)=+\infty.
\end{equation}

\begin{lemma}[One-shell capacity]\label{cost:lem:shell}
In setting \textup{(LT)} let $h\ge1$.
\begin{enumerate}[label=\textup{(\roman*)}]
\item Let $0\le m<e_h$ and let $\alpha_1,\ldots,\alpha_m\in\Cp$ have $v_F(\alpha_i)=r_h$,
repetitions allowed. Then some $z\in Z_h$ satisfies
\[
 \sum_{i=1}^m\bigl(v_F(z-\alpha_i)-r_h\bigr)\le U_h(m).
\]
\item Let $1\le k<h$ and let $A\subset(\OO_K/\pi^h)^\times$ be a set of representatives of the
unit classes modulo $\pi^k$. Then every $z\in Z_h$ satisfies
\[
 \sum_{a\in A}\bigl(v_F(z-[a]t_h)-r_h\bigr)\ge\frac{k\,e_k}{e_h}=U_h(e_k),
\]
with equality for some $z$.
\item In setting \textup{(C)} the bound in \textup{(i)} is attained for every $m<e_h$:
if $\xi$ is a primitive $p^h$-th root of unity and $\alpha_i=\xi^{b_i}-1$, where
$b_1<\cdots<b_m$ are the first $m$ positive integers prime to $p$, then every $z\in Z_h$ has
$\sum_i(v(z-\alpha_i)-r_h)\ge U_h(m)$.
\end{enumerate}
With $m=e_h$, one root at every point of $Z_h$ makes the excess infinite everywhere.
\end{lemma}

\begin{proof}
(i) For each $\alpha_i$ select a unit residue $b_i$ maximizing $v_F([b_i]t_h-\alpha_i)$. For
every unit residue $a$,
\[
 v_F([a]t_h-\alpha_i)\le v_F([a]t_h-[b_i]t_h).
\]
If $a=b_i$ the right side is infinite. Otherwise a strict reverse inequality would force
$v_F([b_i]t_h-\alpha_i)<v_F([a]t_h-\alpha_i)$ by the ultrametric inequality, contradicting
maximality. Thus projecting each root to a nearest point of the shell can only increase the
sum at every point, and we may reason about a multiset of $m$ marked leaves $b_i$. Choose a
least-occupied first branch of the tree and then a least-occupied child at every stage. At
depth $j$ the chosen class contains at most $\floor{m/e_j}$ marks; this follows inductively
from $\floor{\floor{m/e_j}/q}=\floor{m/e_{j+1}}$. Since $m<e_h$, the final leaf $a$ is
unmarked. By \eqref{cost:eq:distance}, summing first over $i$ and then over depths counts
the marks in the selected classes:
\[
 \sum_i\bigl(v_F([a]t_h-[b_i]t_h)-r_h\bigr)=\frac1{e_h}\sum_{j=1}^{h-1}e_j\,N_j(a)\le U_h(m),
\]
where $N_j(a)$ is the number of marks congruent to $a$ modulo $\pi^j$.

(ii) For $j\le k$ exactly $e_k/e_j$ elements of $A$ are congruent to a given unit modulo
$\pi^j$. If $z=[b]t_h$ with $b\notin A$, the same identity gives a sum of at least
$e_h^{-1}\sum_{j\le k}e_j\cdot e_k/e_j=ke_k/e_h$, with equality when the element of $A$
congruent to $b$ modulo $\pi^k$ differs from $b$ modulo $\pi^{k+1}$. If $b\in A$ the sum is
infinite. Finally $U_h(e_k)=ke_k/e_h$ because $\floor{e_k/e_j}=0$ for $j>k$.

(iii) In each block of $e_j$ consecutive positive integers prime to $p$ every unit residue
modulo $p^j$ appears once. Hence every unit class modulo $p^j$ contains at least
$\floor{m/e_j}$ marks. An unmarked point has excess at least $U_h(m)$ by the identity above,
and a marked point has infinite excess.
\end{proof}

The nearest-point projection is why the lemma covers arbitrary roots in $\Cp$ and not only
torsion points: neither algebraicity nor Galois stability of the multiset is assumed. In
terms of Lemma~\ref{cost:lem:greedy}, applied to $X=Z_h$ and $\omega=r_h$, the lemma says
\begin{equation}\label{cost:eq:EU}
 E_{Z_h}(m)\le U_h(m)\quad(0\le m\le e_h),\qquad
 E_{Z_h}(e_k)\ge U_h(e_k)\quad(1\le k<h),
\end{equation}
with $E_{Z_h}=U_h$ in setting \textup{(C)}. Indeed a root of a monic polynomial off the
radius $r_h$ contributes at most $r_h$ at every point of the shell, and
Lemma~\ref{cost:lem:shell}(i) bounds the contribution of the roots on the radius. For the same
reason, when Theorem~\ref{cost:thm:finite} is applied to a union of shells with
$\rho\ge\max r_h$, the roots counted on the part $Z_h$ in
Theorem~\ref{cost:thm:finite}(a) may be taken of valuation exactly $r_h$: a root of larger
valuation and weight $r_h=\rho$ contributes nothing at the points of $Z_h$. Exhaustive
enumerations confirm the value of the maximin in small shells; the first of the following
files covers $14{,}872$ layouts in eight shells, including repeated roots, empty layouts
and full shells.
\costfiles{
\item \artifact{grh-torsion-and-zero-data/verification/check_minimax.txt}
\item \artifact{grh_sharp_transfers_2026-11-25/data/shell_exhaustive_checks.json}}

Combining Theorem~\ref{cost:thm:finite} with Lemma~\ref{cost:lem:shell} gives the aggregate cost of any
finite set of cyclotomic shells at any tolerance.

\begin{corollary}[Any finite set of shells, any tolerance; \stU]\label{cost:cor:shells}
In setting \textup{(C)} let $S$ be a nonempty finite set of positive integers,
$\rho\ge\max_{h\in S}r_h$ and $\delta\ge0$. In increasing order of $h\in S$ define
\begin{equation}\label{cost:eq:shellrec}
 m_h=\min\Bigl\{0\le m\le e_h:\ U_h(m)>\delta+\sum_{s\in S,\ s<h}m_s(r_s-r_h)\Bigr\}.
\end{equation}
Then the least value of $G_\rho(g)$ among integral polynomials $g$ blocking
$\bigcup_{h\in S}Z_h$ at tolerance $\delta$ is $\sum_{h\in S}m_hr_h$, and the least degree is
$\sum_{h\in S}m_h$. Both are attained by the monic polynomial with roots $\xi_h^b-1$, where
$\xi_h$ is a primitive $p^h$-th root of unity and $b$ runs through the first $m_h$ positive
integers prime to $p$; its coefficients lie in $\Qp(\mu_{p^{\max S}})$. Every blocker has at
least $m_h$ roots of valuation $r_h$ for each $h\in S$. The budget threshold also holds for
integral power series.
\end{corollary}

\begin{proof}
The parts of $\Sigma=\bigcup_{h\in S}Z_h$ are the shells, with weights $r_h$, and decreasing
weight is increasing $h$. Since $E_{Z_h}=U_h$ by \eqref{cost:eq:EU}, the recursion
\eqref{cost:eq:rec} is \eqref{cost:eq:shellrec}. By Lemma~\ref{cost:lem:shell}(iii) the first $m_h$ unit exponents form
an admissible initial segment of extremal points on each shell, and the proof of
Theorem~\ref{cost:thm:finite}(b) applies to them.
\end{proof}

\subsection{The exact minimax theorem for consecutive shells}\label{cost:sec:minimax}
{\raggedright\sources{grh_sharp_transfers_2026-11-25; grh_method_transfers_sharp_torsion; grh_sharp_torsion_and_symplectic_blocks; grh_exact_torsion_transfer_2026-11-25}}

For consecutive shells and tolerance zero the recursion has a closed solution, and the same
closed formula holds for formal groups of every finite height. Besides the settings (C) and
(LT) of Definition~\ref{cost:def:lt} we use the following one.

\begin{definition}[Finite height over an unramified base]\label{cost:def:height}
Setting \textup{(H)}: $K/\Qp$ is a finite unramified extension, $F$ is a one-dimensional
commutative formal group law over $\OO_K$ whose reduction has finite height $d$, a fixed
integral formal coordinate is used, $q=p^d$, $v_F=v$, and
\[
 e_h=(q-1)q^{h-1},\qquad r_h=e_h^{-1},\qquad
 Z_h=\set{x:[p^h]x=0,\ [p^{h-1}]x\ne0}\quad(h\ge1),
\]
where $x$ ranges over the maximal ideal of $\OO$. Setting \textup{(C)} is the case
$F=\widehat{\mathbb G}_m$, $d=1$.
\end{definition}

In the three settings write $\varpi=\pi$ in (LT) and $\varpi=p$ in (C) and (H), so that $Z_h$
is the set of points of exact order $\varpi^h$ and $F[\varpi^k]$ is the group of points
killed by $\varpi^k$. Gauss valuations and the notions of Definition~\ref{cost:def:aggregate} are
formed with $v_F$.

\begin{lemma}[The torsion metric at finite height]\label{cost:lem:height}
In setting \textup{(H)}:
\begin{enumerate}[label=\textup{(\roman*)}]
\item every point of $Z_h$ has valuation $r_h$;
\item $F[p^h]$ has $q^h$ points, $Z_h$ has $e_h$ points, and
$[p^k]\colon F[p^H]\to F[p^{H-k}]$ is onto for $0\le k\le H$;
\item $v(x-y)=v(x\ominus y)$ for points of positive valuation, and an integral change of
formal coordinate with unit linear coefficient preserves these distances.
\end{enumerate}
\end{lemma}

\begin{proof}
In $[p](T)$ the coefficient of $T$ is $p$, the coefficients of degrees $2,\ldots,q-1$ belong
to $p\OO_K$, and the coefficient of $T^q$ is a unit; the last two facts follow from the
definition of height and from the unramified base. For $t=v(z)>0$ the unique term of least
valuation in $[p](z)$ has valuation $qt$ if $t<r_1$ and $1+t$ if $t>r_1$. In the latter
region repeated application of $[p]$ increases the valuation by one and never gives zero. At
$t=r_1$ the image has valuation at least $1+r_1$. Thus a nonzero point killed by $p$ has
valuation exactly $r_1$, and induction shows that a point of exact order $p^h$ has valuation
$r_h$. The first unit coefficient of $[p^h](T)$ occurs in degree $q^h$, by composition after
reduction. Weierstrass preparation gives $q^h$ roots in the open unit disc, counted with
multiplicity. They are simple, since translation by a torsion point identifies the local
zero with the zero at the origin, where the derivative is $p^h\ne0$. The count of $Z_h$ and
the surjectivity of $[p^k]$ follow from the orders of the groups. For (iii), the series
$F(X,[-1]Y)$ is $(X-Y)$ times a series with unit constant term. For elliptic formal groups
see also~\cite[\S4, Lemma~4.1]{BoxallGrantSingular}.
\end{proof}

\begin{theorem}[Exact minimax cost of blocking torsion specialization; \stU]\label{cost:thm:minimax}
Consider one of the settings \textup{(C)}, \textup{(LT)}, \textup{(H)}. Let $h_0\ge1$,
$H\ge h_0$, $N=H-h_0+1$ and $\rho\ge r_{h_0}$.
\begin{enumerate}[label=\textup{(\roman*)}]
\item The minimum of $\Delta_\rho(g)$ among all nonzero polynomials $g\in\Cp[T]$ that block
the shells $Z_{h_0},\ldots,Z_H$ at tolerance zero is
\[
 N\,q^{1-h_0}.
\]
It is attained by a monic polynomial with coefficients in $\OO$, of degree $q^N-1$, whose
roots lie on these shells.
\item In settings \textup{(C)} and \textup{(LT)}, every such blocker has at least
$(q-1)q^{h-h_0}$ roots of valuation $r_h$ for each $h_0\le h\le H$. In particular its degree
is at least $q^N-1$.
\item In settings \textup{(C)} and \textup{(H)}, if $\Delta_\rho(g)<Nq^{1-h_0}$, then for
some $h_0\le h\le H$ an entire coset of $F[\varpi^{h_0-1}]$ contained in $Z_h$ consists of
good points. In setting \textup{(C)} this is a set
$\set{\zeta\eta-1:\eta\in\mu_{p^{h_0-1}}}$ with $\zeta$ of order $p^h$. If $\deg g<q^N-1$,
then $g$ has a good point on these shells whatever the value of $\Delta_\rho(g)$ is.
\end{enumerate}
\end{theorem}

The coefficient field is unrestricted, and this is part of the sharpness statement: the
extremal polynomials have coefficients in a finite extension whose ramification grows with
$H$. A restriction to a fixed coefficient field imposes Galois symmetry on the roots and
changes the answer (\cref{cost:sec:fields}).

\begin{proof}[Proof of \textup{(iii)}, and of \textup{(i)} in settings \textup{(C)} and \textup{(H)}]
Write $k=h_0-1$, $M=q^k$ and $\Lambda=F[\varpi^N]$. By Lemma~\ref{cost:lem:height} the map
$[\varpi^k]\colon F[\varpi^H]\to\Lambda$ is onto, and the points of the shells
$Z_{h_0},\ldots,Z_H$ form the inverse image of $\Lambda\setminus\set0$. Choose one lift $x_u$
of each nonzero $u\in\Lambda$. A lift of a point of exact order $\varpi^j$ has exact order
$\varpi^{j+k}$, and the difference $x_u\ominus x_w$ is a lift of $u\ominus w$. Hence, for
distinct nonzero $u,w$,
\[
 v(x_u)=v(u)/M,\qquad v(x_u-x_w)=v(u\ominus w)/M.
\]
Each shell in $\Lambda$ has total valuation $e_jr_j=1$. Let $P(T)=\prod_{u\ne0}(T-x_u)$,
a monic integral polynomial of degree $q^N-1$. Since every $v(x_u)\le r_{h_0}\le\rho$,
Lemma~\ref{cost:lem:roots} gives $G_\rho(P)=N/M$ and $G_0(P)=0$. Translation permutes $\Lambda$:
the differences $u\ominus w$, for $w\ne0,u$, range over all nonzero elements of $\Lambda$
except $u$. Therefore
\[
 v(P'(x_u))=\frac{N-v(u)}M=\frac NM-v(x_u),
\]
and every Lagrange basis polynomial $\ell_u(T)=P(T)/((T-x_u)P'(x_u))$ has $G_\rho(\ell_u)=0$.

Let $g$ be a nonzero polynomial with $\Delta_\rho(g)<N/M$, normalized so that $G_0(g)=0$,
and put $A=G_\rho(g)$. Divide: $g=PQ+R$ with $\deg R<q^N-1$. Monic division over $\OO$
leaves $Q$ and $R$ integral, so $G_\rho(PQ)\ge N/M>A$ and $G_\rho(R)=A$. Interpolation gives
\[
 R(T)=\sum_{u\ne0}g(x_u)\,\ell_u(T),\qquad A\ge\min_{u\ne0}v(g(x_u)).
\]
This holds for every choice of the lifts. If each of the $q^N-1$ fibres of $[\varpi^k]$ over
$\Lambda\setminus\set0$ contained a point that is not good, we could select these points as
lifts and contradict the inequality. Hence one fibre consists of good points. A fibre is a
coset of $F[\varpi^k]$, and its points have a common order. If instead $\deg g<q^N-1$, then
$R=g$ and the same interpolation applies whatever $A$ is. This proves (iii) and the lower
bound in (i).

For the attainment in (i), take $g=P$. At a selected node its value is zero. Let $z$ be
another point of the fibre of $x_u$; this occurs only for $k>0$. Then $z\ominus x_u$ is a
nonzero element of $F[\varpi^k]$ and has valuation at least $r_k>r_{k+1}$, which is larger
than all the valuations $v(x_u-x_w)$ and than $v(x_u)$. The remaining factors therefore
have the same valuations at $z$ as at $x_u$, and
\[
 v(P(z))=\frac NM-v(x_u)+v(z-x_u)>\frac NM=G_\rho(P).
\]
Thus $P$ blocks all $N$ shells at cost exactly $Nq^{1-h_0}$.
\end{proof}

\begin{proof}[Proof of \textup{(ii)}, and of \textup{(i)} in settings \textup{(C)} and \textup{(LT)}]
In setting \textup{(C)} this is the argument of the source. In setting \textup{(LT)} the
source states that the cyclotomic argument carries over unchanged with $q$ in place of $p$:
it uses only that $Z_h$ has $e_h$ points, all of valuation $r_h$, and that the unit
residues modulo $\pi^h$ form a tree with $q-1$ first branches and $q$ children at every
later vertex, with the distances \eqref{cost:eq:distance}. These facts are
Lemma~\ref{cost:lem:metric}, and they are exactly what Lemma~\ref{cost:lem:shell} uses.
We write the argument once, with $q$ and $v_F$.

Let $g$ be a nonzero polynomial blocking $Z_{h_0},\ldots,Z_H$ at tolerance zero, write
$c=q^{1-h_0}$, and let $m_h$ be the number of roots of $g$ of valuation $r_h$. For
$h_0\le h\le H$ the slack $s_h=G_\rho(g)-G_{r_h}(g)$ equals, by \eqref{cost:eq:gauss},
$\sum_\alpha\bigl(\min\set{\rho,v_F(\alpha)}-\min\set{r_h,v_F(\alpha)}\bigr)$. Every
summand is nonnegative, and a root of valuation $r_k$ with $h_0\le k<h$ contributes
$r_k-r_h$; hence
\begin{equation}\label{cost:eq:slack}
 s_h\ge\sum_{k=h_0}^{h-1}(r_k-r_h)\,m_k ,
\end{equation}
with equality if every root of $g$ lies on one of the shells $h_0,\ldots,H$. By
Lemma~\ref{cost:lem:roots}(ii), for $z\in Z_h$,
\begin{equation}\label{cost:eq:shellslack}
 v_F(g(z))-G_\rho(g)=\sum_{v_F(\alpha)=r_h}\bigl(v_F(z-\alpha)-r_h\bigr)-s_h .
\end{equation}

We show by induction on $h$ that $m_h\ge M_h:=(q-1)q^{h-h_0}$. Put $k=h-h_0$, so that
$M_h=e_{k+1}$ and $M_h/e_h=c$. By \eqref{cost:eq:slack} and the induction hypothesis,
\begin{equation}\label{cost:eq:slackL}
 s_h\ge\sum_{i=h_0}^{h-1}(r_i-r_h)M_i=kc-\frac{q^k-1}{e_h}=\frac{kM_h-(q^k-1)}{e_h}.
\end{equation}
The last expression equals $U_h(M_h-1)$. Indeed $e_j$ divides $M_h$ for $j\le k+1$, and
$e_j>M_h-1$ for $j>k+1$, so
\[
 e_hU_h(M_h-1)=\sum_{j=1}^{k+1}(M_h-e_j)=(k+1)M_h-(q^{k+1}-1)=kM_h-(q^k-1);
\]
when $k+1=h$ the term $j=h$, which does not occur in $U_h$, is zero, so this also covers
$h_0=1$. For $k=0$ both sides of \eqref{cost:eq:slackL} are zero, which is the base of the
induction. Suppose $m_h<M_h$. Then $m_h<e_h$, and Lemma~\ref{cost:lem:shell}(i), the
monotonicity of $U_h$ on $\set{0,\ldots,e_h-1}$ and \eqref{cost:eq:slackL} give a point
$z\in Z_h$ at which the sum in \eqref{cost:eq:shellslack} is at most
$U_h(m_h)\le U_h(M_h-1)\le s_h$. Then $v_F(g(z))\le G_\rho(g)$, contrary to blocking.
Hence $m_h\ge M_h$ for every $h_0\le h\le H$. This is (ii), and summing gives
$\deg g\ge\sum_{k=1}^N(q-1)q^{k-1}=q^N-1$. By \eqref{cost:eq:budget},
$\Delta_\rho(g)\ge\sum_hm_hr_h\ge Nc$, the lower bound in (i).

For the attainment take on each shell $Z_h$ the points $[a]t_h$, with $a$ running through
a set of representatives of the unit classes modulo $\pi^{k+1}$, $k=h-h_0$; in setting
\textup{(C)} these are the points $\xi_h^a-1$ with $1\le a\le p^{k+1}$, $p\nmid a$. Let $g$
be the monic polynomial with these $\sum_hM_h=q^N-1$ roots. All roots have positive
valuation at most $r_{h_0}\le\rho$, so $G_0(g)=0$ and $G_\rho(g)=\sum_hM_hr_h=Nc$. At a
point of $Z_h$ the slack is exactly the middle term of \eqref{cost:eq:slackL}, by the
equality case of \eqref{cost:eq:slack}. If $k+1=h$, that is $h_0=1$, the roots exhaust
$Z_h$ and $g$ vanishes there. Otherwise Lemma~\ref{cost:lem:shell}(ii), applied with
$k+1<h$, shows that the sum in \eqref{cost:eq:shellslack} is at least
$U_h(e_{k+1})=(k+1)c$ at every point of $Z_h$, which exceeds the slack by
$(q^{k+1}-1)/e_h>0$. By \eqref{cost:eq:shellslack} the polynomial blocks every shell.
\end{proof}

The second proof shows where the value comes from: the slack
$G_\rho(g)-G_{r_h}(g)$ at shell $h$ contains the savings supplied by the roots on all
earlier shells, and keeping these savings in full, instead of discarding them through a
positive part, is what distinguishes the aggregate convention from the packet convention of
\cref{cost:sec:packet}.

\begin{corollary}[Optimality of the universal endpoint and the degree trade-off; \stU]\label{cost:cor:endpoint}
In the settings of Theorem~\ref{cost:thm:minimax} let $B\ge0$ and
\[
 H_*(B)=h_0+\floor{B\,q^{h_0-1}}.
\]
Every nonzero polynomial $g$ with $\Delta_\rho(g)\le B$, in particular every nonzero
integral $g$ with $G_\rho(g)\le B$, has a good point at tolerance zero on some shell $Z_h$
with $h_0\le h\le H_*(B)$; in settings \textup{(C)} and \textup{(LT)} this is
\cref{tors:thm:selection}\textup{(a)} and \cref{tors:cor:lt}. No smaller endpoint has this
property. In settings \textup{(C)} and \textup{(H)}, if in addition $\deg g\le D_0$ is
imposed, the least endpoint is
\[
 h_0+\min\bigl\{\floor{q^{h_0-1}B},\ \floor{\log_q(D_0+1)}\bigr\}.
\]
\end{corollary}

\begin{proof}
By Theorem~\ref{cost:thm:minimax}(i), a window of $N$ shells can be blocked by a polynomial
with $\Delta_\rho\le B$ exactly when $Nq^{1-h_0}\le B$. For $N=\floor{Bq^{h_0-1}}+1$ this
is the sufficiency of $H_*(B)$, which in setting \textup{(H)} is not covered by
\cref{sec:tors}. If $\floor{Bq^{h_0-1}}\ge1$, the extremal polynomial for the window of
$\floor{Bq^{h_0-1}}$ shells shows that $H_*(B)-1$ does not suffice; if it is zero, the
first allowed shell is the least possible endpoint. With a degree bound, a window of $N$
shells works for every admissible $g$ precisely when $N>q^{h_0-1}B$ or $q^N-1>D_0$: below
either threshold Theorem~\ref{cost:thm:minimax}(iii) applies, and when both inequalities
fail the extremal polynomial of Theorem~\ref{cost:thm:minimax}(i) meets both constraints.
\end{proof}

The sufficiency of $H_*(B)$ for power series in several variables, with independent
choices of the coordinates, and for the sum of the valuations of a finite packet, is
\cref{tors:thm:selection}\textup{(b)} and \cref{tors:cor:packets}; a polynomial in one
coordinate already forces the stated endpoint. In setting (C) the corollary determines the
optimal order exponent left open by the earlier sufficient bound
$h_0+\ceil{p^{h_0}B(h_0+B+2)}$ for exact specialization (\cref{tors:tab:endpoints}).

\begin{example}\label{cost:ex:minimax}
For $p=3$, $h_0=2$, $B=1$ and any $\rho\ge1/6$, the optimal endpoint is $H_*=5$. The extremal polynomial for the window $h=2,3,4$ has
$2$, $6$ and $18$ roots on these shells, degree $26$ and budget exactly $1$; no polynomial
of degree at most $25$ blocks this window at any budget. Extending the window to $H=5$
guarantees three good roots forming a coset of $\mu_3$. In setting (H), the reduction of
$y^2=x^3-x$ at $3$ is supersingular, so its formal group has height two and $q=9$; with
$h_0=2$, $B=1$ and $\rho\ge1/72$ the least endpoint is $H_*=11$, and at $H=10$ the extremal
polynomial uses one lift of each of the $9^9-1$ nonzero points of $F[3^9]$. Exact checks of
the cyclotomic statements in small cases, and of small cases at heights two and three (last
file), are recorded in:
\costfiles{
\item \artifact{grh_sharp_transfers_2026-11-25/data/aggregate_minimax_checks.json}
\item \artifact{grh_sharp_torsion_and_symplectic_blocks/checks/results.txt}
\item \artifact{grh_method_transfers_sharp_torsion/checks/torsion_checks.txt}
\item \artifact{grh_method_transfers_sharp_torsion/checks/formal_height_checks.txt}}
\end{example}

\begin{remark}[Fourier coefficients of measures]\label{cost:rem:measures}
Let $\mu$ be a bounded $\OO$-valued measure on $\Zp$ with Mahler moments
$a_j=\int\binom xj\,\dd\mu(x)$. For a root of unity $\zeta$ of $p$-power order,
$\widehat\mu(\zeta)=\int\zeta^x\,\dd\mu(x)=\sum_{j\ge0}a_j(\zeta-1)^j$, by uniform convergence
of the binomial expansion on $\Zp$. Thus a single moment with $v(a_j)+j\rho\le B$ guarantees
an entire coset of characters of conductor $p^h$, $h_0\le h\le H_*(B)$, each with Fourier
coefficient of valuation at most $\min_i(v(a_i)+i\rho)$. Sharpness persists among finite
atomic measures supported on nonnegative integers: each $T^j=((1+T)-1)^j$ is an integral
combination of the transforms $(1+T)^r$, $0\le r\le j$, of point masses, so the extremal
polynomial of Theorem~\ref{cost:thm:minimax} is realized with support in $\set{0,\ldots,p^N-1}$.
\end{remark}

\subsection{The optimal exponent for the fixed-radius comparison}\label{tors:sec:optimal}
\sources{grh_sharp_transfers_2026-11-25; grh_exact_torsion_transfer_2026-11-25}

Fix $h_0\ge1$ and a comparison radius $\rho\ge r_{h_0}$. A nonzero polynomial
$g\in\Cp[T]$ \emph{blocks} the shells $h_0,\dots,H$ if
\[
  v(g(z))>G_\rho(g)\qquad\text{for every }z\in\torsZ{h_0}{H}.
\]
The radius $\rho$ is fixed throughout; it is not replaced by the radius $r_h$ of the
point that is tested. This is the aggregate convention of \cref{cost:def:aggregate} at tolerance zero, for the
set $\torsZ{h_0}{H}$. Its exact price is determined in \cref{sec:cost}: by
\cref{cost:thm:minimax}\textup{(i)} every nonzero $g\in\Cp[T]$ that blocks the shells
$h_0,\dots,H$ satisfies
\begin{equation}\label{tors:eq:budget}
  \Delta_\rho(g)\ \ge\ (H-h_0+1)\,p^{1-h_0},
\end{equation}
and by \cref{cost:thm:minimax}\textup{(ii)} it has at least $(p-1)p^{h-h_0}$ roots of
valuation $r_h$ on each shell $h_0\le h\le H$. The proof combines the one-shell capacity
\cref{cost:lem:shell} with the saving \eqref{tors:eq:slack} from earlier shells.

\begin{theorem}[Optimal torsion selection; \stU]\label{tors:thm:selection}
Let $h_0\ge1$, $\rho\ge r_{h_0}$ and $B\ge0$, and put
\[
  \torsHs(B)=h_0+\floor[\big]{B\,p^{h_0-1}} .
\]
\begin{enumerate}[label=\textup{(\alph*)}]
\item If $g\in\Cp[T]$ is nonzero and $\Delta_\rho(g)\le B$, there are $h$ with
  $h_0\le h\le\torsHs(B)$ and $z\in Z_h$ such that
  \[
     v(g(z))\le G_\rho(g).
  \]
\item If $F$ is a nonzero bounded series in $s\ge1$ variables and $\Delta_\rho(F)\le B$,
  there is a torsion tuple $z\in\torsZ{h_0}{\torsHs(B)}^s$ with
  \[
     v(F(z))\le G_\rho(F).
  \]
\end{enumerate}
In particular both conclusions hold for integral $g$ and $F$ with $G_\rho\le B$, and for
integral $F$ having a witness $x$ of depth $\rho$ with $v(F(x))\le B$; in the latter case
$v(F(z))\le v(F(x))$. The endpoint $\torsHs(B)$ does not depend on $s$, on degrees or on
the field generated by the coefficients.
\end{theorem}

\begin{proof}
(a) Put $N=\floor{Bp^{h_0-1}}+1$. If $g$ blocked the $N$ shells $h_0,\dots,\torsHs(B)$,
\eqref{tors:eq:budget}, that is \cref{cost:thm:minimax}\textup{(i)}, would give
$\Delta_\rho(g)\ge Np^{1-h_0}>B$.

(b) First let $F$ be a polynomial, and put $A=G_\rho(F)$, $a=G_0(F)$. Write
\[
  F=\sum_\beta f_\beta(T_1)\,T_2^{\beta_2}\cdots T_s^{\beta_s}
\]
and choose $\beta$ attaining $A$, so that $G_\rho(f_\beta)+\rho\abs\beta=A$. Every
coefficient of $f_\beta$ has valuation at least $a$, hence
$\Delta_\rho(f_\beta)\le G_\rho(f_\beta)-a\le A-a\le B$. By (a) there is
$z_1\in\torsZ{h_0}{\torsHs(B)}$ with $v(f_\beta(z_1))\le G_\rho(f_\beta)$. The polynomial
$F(z_1,T_2,\dots,T_s)$ has the coefficient $f_\beta(z_1)$ at $T^\beta$, so its Gauss
valuation is at most $A$; in particular it is nonzero. All its coefficients have
valuation at least $a$, because $v(z_1)>0$. Its variation $\Delta_\rho$ is therefore at
most $A-a\le B$, and the argument can be repeated. After $s$ steps one obtains a constant
of valuation at most $A$.

For a bounded series keep the real number $a=G_0(F)$ as a lower bound for the coefficient
valuations; no normalization by a scalar of valuation $a$ is used, since such a scalar
need not exist. Put $H=\torsHs(B)$ and $A=G_\rho(F)$, and apply
\cref{tors:lem:truncation} with $M=A$ and $r=r_H\le\rho$. The truncation $P$ satisfies
$G_\rho(P)=A$ and $G_0(P)\ge a$, so $\Delta_\rho(P)\le B$. The polynomial case gives
$z\in\torsZ{h_0}{H}^s$ with $v(P(z))\le A$, and then $v(F(z))=v(P(z))\le A$ by
\cref{tors:lem:truncation}(a).

The last assertions follow from $\Delta_\rho\le G_\rho$ for integral series and from
\eqref{tors:eq:witness}.
\end{proof}

\begin{remark}[Sharpness, and the bounds that are recovered]\label{tors:rem:recovered}
(1) By \cref{cost:thm:minimax}\textup{(i)} the bound \eqref{tors:eq:budget} is attained
by a monic integral polynomial of degree $p^{H-h_0+1}-1$ with roots in the blocked shells.
Hence $\torsHs(B)$ is the least endpoint valid for all $g$, uniformly in the degree and in
the coefficient field (\cref{cost:cor:endpoint}). A restriction to a fixed finite extension of $\Qp$ imposes Galois
symmetry on the roots and can lower the endpoint; see \cref{tors:sec:fixedfield} and
\cref{sec:cost}.

(2) \Cref{tors:thm:selection} has the hypotheses and the conclusion of the earlier
zero-tolerance bound with endpoint $h_0+\ceil{p^{h_0}B(h_0+B+2)}$, which was obtained from
a recurrence for the root masses on successive shells, and it replaces that bound for one
series and for finite families; the earlier endpoint is kept only in
\cref{tors:cor:packets}\textup{(b)}, where it is the endpoint of the source statement.
\Cref{tors:tab:endpoints} compares the two endpoints.

(3) If $\Delta_\rho(g)<r_{h_0}$, then \emph{every} $z\in Z_{h_0}$ satisfies
$v(g(z))\le G_\rho(g)$: by \cref{tors:lem:shells}(a) no root of $g$ has valuation
$r_{h_0}$, and \eqref{tors:eq:value} gives $v(g(z))=G_{r_{h_0}}(g)\le G_\rho(g)$.

(4) The endpoint bounds the exponent, so the orders of the selected roots of unity are at
most $p^{\torsHs(B)}$; it is not asserted that the largest order occurs. The grid
$\torsZ{h_0}{\torsHs(B)}^s$ has $(p^{\torsHs(B)}-p^{h_0-1})^s$ points. Independence of the
exponent from $s$ does not make a search over the grid independent of $s$.
\end{remark}

\begin{table}[ht]
\centering\small
\begin{tabular}{@{}rrrrr@{}}
\toprule
$p$ & $h_0$ & $B$ & $h_0+\ceil{p^{h_0}B(h_0+B+2)}$ & $\torsHs(B)=h_0+\floor{Bp^{h_0-1}}$\\
\midrule
2&1&1&9&2\\
2&1&10&261&11\\
3&2&1&47&5\\
3&2&5&407&17\\
2&3&10&1203&43\\
5&2&2&302&12\\
\bottomrule
\end{tabular}
\caption{The earlier sufficient endpoint for the exact fixed-radius comparison and the
optimal endpoint of \cref{tors:thm:selection}, for the same data
(\artifact{grh_sharp_transfers_2026-11-25/data/bound_comparison.csv}).}
\label{tors:tab:endpoints}
\end{table}

\begin{corollary}[Aggregate comparison for families; \stU]\label{tors:cor:packets}
Let $h_0,\rho,B$ be as in \cref{tors:thm:selection}.
\begin{enumerate}[label=\textup{(\alph*)}]
\item Let $F_1,\dots,F_n$ be nonzero bounded series in the same $s$ variables with
  $\sum_i\Delta_\rho(F_i)\le B$. There is $z\in\torsZ{h_0}{\torsHs(B)}^s$ with
  \[
    \sum_{i=1}^nv(F_i(z))\le\sum_{i=1}^nG_\rho(F_i).
  \]
  For integral $F_i$ it suffices that $\sum_iG_\rho(F_i)\le B$, or that there are
  witnesses $x_i$ of depth $\rho$, possibly different for different $i$, with
  $\sum_iv(F_i(x_i))\le B$. In the integral case every $F_i(z)$ is nonzero and
  $v(F_i(z))\le B$.
\item Let $(F_i)_{i\in I}$ be an arbitrary family of nonzero integral series in $s$
  variables such that $\sum_{i\in J}G_\rho(F_i)\le B$ for every finite $J\subseteq I$, and
  put $H_0=h_0+\ceil{p^{h_0}B(h_0+B+2)}$. There is a single $z$ in the grid
  $\torsZ{h_0}{H_0}^s$ such that $\sum_{i\in J}v(F_i(z))\le B$ for every finite
  $J\subseteq I$.
\end{enumerate}
\end{corollary}

\begin{proof}
(a) Apply \cref{tors:thm:selection}(b) to the nonzero bounded series $\prod_iF_i$ and use
the additivity of $G_\rho$ and $\Delta_\rho$ (\cref{tors:lem:gauss}(a)). For integral
series $\Delta_\rho(F_i)\le G_\rho(F_i)\le v(F_i(x_i))$, and all values $F_i(z)$ are
integral, so a finite bound for the sum of their valuations bounds each of them.

(b) For a finite $J\subseteq I$ let $S_J$ be the set of $z$ in the grid
$\torsZ{h_0}{H_0}^s$ with $\sum_{i\in J}v(F_i(z))\le B$. Since $\torsHs(B)\le H_0$, part (a)
shows that every $S_J$ is nonempty, and
$S_{J_1\cup\dots\cup J_k}\subseteq S_{J_1}\cap\dots\cap S_{J_k}$ because the valuations
are nonnegative. The sets $S_J$ are subsets of one finite set and have the finite
intersection property, so their intersection is nonempty. This is a compactness argument
on the finite grid, not a statement about an infinite product.
\end{proof}

The conclusion is aggregate. It does not assert $v(F_i(z))\le G_\rho(F_i)$ for each $i$,
and by the examples of \cref{sec:cost} no such assertion with an exponent depending only
on $p$, $h_0$ and $B$ is true over $\OO$. This matters for applications that impose
separate bounds on individual Selmer groups; applications that charge a sum of
denominator lengths can use \cref{tors:cor:packets} directly.

\paragraph{Lubin--Tate torsion.}
The passage from one variable to several variables and to bounded series in
\cref{tors:thm:selection}(b) uses only Gauss valuations and truncation, and the
one-variable statement holds for Lubin--Tate torsion by \cref{cost:thm:minimax}. Let
$K/\Qp$ be finite with uniformizer $\pi$ and residue field of
cardinality $q$, let $w$ be the valuation of $\Cp$ with $w(\pi)=1$, and let $\mathcal F$ be
a one-dimensional Lubin--Tate formal $\OO_K$-module for $\pi$ \cite{LubinTate1965Formal}.
Write $Z_h^{\mathcal F}$ for the set of points of exact $\pi^h$-torsion of $\mathcal F$.
Here the formal parameter itself is evaluated; nothing is subtracted.

\begin{corollary}[Lubin--Tate selection; \stU]\label{tors:cor:lt}
Put $e_h=(q-1)q^{h-1}$ and $r_h=1/e_h$, and define Gauss valuations with $w$ in place of
$v$. Let $h_0\ge1$, $\rho\ge r_{h_0}$ and $B\ge0$. For every nonzero bounded series $F$ in
$s$ variables with $\Delta_\rho(F)\le B$ there are $h_1,\dots,h_s$ in
$[h_0,\,h_0+\floor{Bq^{h_0-1}}]$ and points $z_j\in Z^{\mathcal F}_{h_j}$ with
\[
  w(F(z_1,\dots,z_s))\le G_\rho(F).
\]
\end{corollary}

\begin{proof}
This is setting \textup{(LT)} of \cref{cost:def:lt}, with $v_F=w$. By
\cref{cost:thm:minimax}\textup{(i)} in that setting, a nonzero polynomial blocking the
shells $Z^{\mathcal F}_{h_0},\dots,Z^{\mathcal F}_H$ has
$\Delta_\rho\ge(H-h_0+1)q^{1-h_0}$. The statement in one variable follows as in
\cref{tors:thm:selection}(a), and the proof of \cref{tors:thm:selection}(b), which uses
only Gauss valuations and truncation, carries over unchanged.
\end{proof}

The endpoint in \cref{tors:cor:lt} is again optimal; the extremal polynomials, and formal
groups of arbitrary finite height, are treated in \cref{sec:cost}.

\subsection{Formal groups over ramified bases and overlapping torsion valuations}\label{cost:sec:ramified}
{\raggedright\sources{grh_method_transfers_sharp_torsion; grh_method_transfers_research}}

The uniform radii $r_h$ were used in Theorem~\ref{cost:thm:minimax} to express the answer through
$q$ and $h_0$. Over a ramified base a formal group of finite height can have torsion points
of one exact order with several valuations, and points of different orders with the same
valuation; cutting by valuation then no longer cuts between two orders, and the closed
formula changes. In this subsection $K/\Qp$ is an arbitrary finite extension, $F$ is a
one-dimensional commutative formal group law over $\OO_K$ whose reduction has finite height
$d$, a fixed integral coordinate is used, $Q=p^d$, and $F[p^j]$ denotes the geometric
$p^j$-torsion in the open unit disc. As in Lemma~\ref{cost:lem:height}, $v(x-y)=v(x\ominus y)$,
and the prepared factorization of $[p^j](T)$ is a unit times $\prod_{y\in F[p^j]}(T-y)$ with
simple roots; since $[p^j]'(0)=p^j$,
\begin{equation}\label{cost:eq:totalval}
 \sum_{0\ne y\in F[p^j]}v(y)=j.
\end{equation}

\begin{proposition}[Valuation cut-off; \stU]\label{cost:prop:cutoff}
Let $\Gamma=F[p^H]$, fix $\rho>0$, and put
\[
 \Gamma_\rho=\set{z\in\Gamma:v(z)>\rho},\qquad
 A_*=\frac{H-\sum_{0\ne z\in\Gamma_\rho}v(z)}{\abs{\Gamma_\rho}}.
\]
Assume $\Gamma\ne\Gamma_\rho$. Among nonzero integral polynomials that block
$\Gamma\setminus\Gamma_\rho$ at tolerance zero, the least value of $G_\rho$ is $A_*$ and the
least degree is $[\Gamma:\Gamma_\rho]-1$; one polynomial attains both. If
$G_\rho(g)<A_*$ or $\deg g<[\Gamma:\Gamma_\rho]-1$, then an entire coset of $\Gamma_\rho$ in
$\Gamma\setminus\Gamma_\rho$ consists of good points.
\end{proposition}

\begin{proof}
The formal group law and inverse preserve the condition $v(z)>\rho$, so $\Gamma_\rho$ is a
subgroup. Adding an element of $\Gamma_\rho$ changes a parameter by an amount of valuation
greater than $\rho$, so the valuation is constant on each nonidentity coset. By
\eqref{cost:eq:totalval} the total valuation of one representative $x_u$ from each
nonidentity coset $u$ is $A_*$. The valuation of the difference of two such representatives
depends only on their difference coset. Let $P$ be the product of the linear factors
$T-x_u$. Then $G_\rho(P)=A_*$, and translation in $\Gamma/\Gamma_\rho$ gives
$v(P'(x_u))=A_*-v(x_u)$, so the Lagrange basis polynomials of these nodes have $G_\rho=0$.
The interpolation argument in the proof of Theorem~\ref{cost:thm:minimax}, with arbitrary
representatives, gives the good coset. Conversely, at another point $z$ of the coset of
$x_u$ all factors of $P$ except $T-x_u$ keep their valuations and
$v(z-x_u)>\rho\ge v(x_u)$, so $v(P(z))>A_*$; thus $P$ blocks.
\end{proof}

The formula needs only the multiset of torsion valuations, which is available from the
Newton polygon of $[p^H]$. If the comparison radius is at least every nonzero torsion
valuation, then $\Gamma_\rho=\set0$ and the obstruction budget is exactly $H$ over every
base; for this it suffices that $\rho\ge\max_{0\ne z\in F[p]}v(z)$, since multiplication by
$p$ strictly increases the valuation of a nonzero point of the disc until it reaches zero.

For a window of exact orders that is not a valuation cut-off, Theorem~\ref{cost:thm:finite} applies
once the metric on the test set is known, and the metric can be read off from torsion
counts. Take $\Gamma=F[p^H]$, $D=F[p^{h_0-1}]$ and the target $\Sigma=\Gamma\setminus D$. The
sets $V_r=\set{x\in\Gamma:v(x)\ge r}$ are subgroups and their cosets are the balls of level
$r$. For each torsion valuation $r$ the numbers $\abs{V_r}$ and $\abs{D\cap V_r}$ are
obtained from the slopes, with multiplicities, of the Weierstrass polynomials of $[p^H]$ and
$[p^{h_0-1}]$. At adjacent levels $r<t$ a coset of $V_r$ has $\abs{V_r}/\abs{V_t}$ children;
if it meets $D$, exactly $\abs{D\cap V_r}/\abs{D\cap V_t}$ of them meet $D$, by translation,
and otherwise none does. Marking the path through the origin, this constructs the finite
tree of balls up to permutation of equivalent children. After the leaves in $D$ are deleted,
the last common ball of two surviving leaves gives the valuation of their difference, and
the comparison with the origin gives each weight. Thus Theorem~\ref{cost:thm:finite} solves the
prescribed order interval from integer populations and rational slopes; there is no need to
solve for the torsion coordinates. For the general theory of the subgroups of torsion points
closest to zero see~\cite{Rabinoff2009Higher}.

\begin{corollary}[Termination; \stU]\label{cost:cor:termination}
Fix $h_0\ge1$, $\rho>0$ and $B\ge0$. For all sufficiently large $H$, every nonzero integral
polynomial or power series $g$ with $G_\rho(g)\le B$ has a good point at tolerance zero in
$F[p^H]\setminus F[p^{h_0-1}]$. Consequently the finite computation above determines the
least such $H$.
\end{corollary}

\begin{proof}
Choose $0<t\le\rho$ smaller than all nonzero valuations in $D=F[p^{h_0-1}]$. The torsion
points of valuation greater than $t$ form a finite subgroup $\Gamma_t$: the series
$[p](X)=pX+O(X^2)$ has integral coefficients and raises a valuation exceeding $t$ by at
least $\min\set{1,t}$ until the point is killed, and a nonzero point of valuation greater
than $1$ is not torsion because from then on the linear term dominates. Hence $\Gamma_t$ is
killed by a bounded power of $p$, and $D\subseteq\Gamma_t\subseteq F[p^H]$ for large $H$.
By Proposition~\ref{cost:prop:cutoff} at radius $t$, if $G_t(g)<A_*=(H-c)/\abs{\Gamma_t}$ with
$c=\sum_{0\ne z\in\Gamma_t}v(z)$, some point $x$ outside $\Gamma_t$, hence outside $D$,
satisfies $v(g(x))\le G_t(g)\le G_\rho(g)$. It remains to take $H$ with $A_*>B$. Power series
are truncated as in Theorem~\ref{cost:thm:finite}(d).
\end{proof}

Deep in the torsion tower the optimal costs scale exactly. Write
$[p](T)=\sum_{i\ge1}a_iT^i$; the coefficient $a_Q$ is a unit and the earlier coefficients
lie in the maximal ideal. Put
\[
 \kappa=\min_{1\le i<Q}\frac{v(a_i)}{Q-i},\qquad\Theta=Q\kappa.
\]
For $z$ in the maximal ideal of $\OO$,
\begin{equation}\label{cost:eq:cap}
 \min\set{v([p]z),\Theta}=Q\min\set{v(z),\kappa}.
\end{equation}
Indeed, for $v(z)<\kappa$ the term $a_Qz^Q$ has strictly smaller valuation than every other
term, and for $v(z)\ge\kappa$ every term has valuation at least $\Theta$. Applied to formal
differences, \eqref{cost:eq:cap} gives
\begin{equation}\label{cost:eq:capdiff}
 \min\set{v(x-y),\kappa}=\frac1Q\min\set{v([p]x-[p]y),\Theta}.
\end{equation}

\begin{proposition}[Exact scaling; \stU]\label{cost:prop:scaling}
Let $\Sigma$ be a finite set of nonzero torsion points of $F$ with all valuations below
$\Theta$, and let $\Sigma'=[p]^{-1}\Sigma$. Let $B$ be the least blocking budget for $\Sigma$
at tolerance $\delta$ and radius $\rho$ in the sense of Theorem~\ref{cost:thm:finite}. If
$B+\delta<\Theta$, then the least blocking budget for $\Sigma'$ at tolerance $\delta/Q$ and
radius $\rho/Q$ is exactly $B/Q$, and the least degree of a budget-minimizing blocker is the
same for $\Sigma$ and $\Sigma'$. When $\rho$ exceeds every valuation of $\Sigma$, both radii
may instead be replaced by one sufficiently large fixed radius.
\end{proposition}

\begin{proof}
By Theorem~\ref{cost:thm:finite} the optima are attained by monic polynomials whose roots are
test points, so it suffices to compare multisets of marked test points. A multiset of marks
of cost $b$ defeats a target $s$ exactly when the sum of its pair valuations against $s$
exceeds $b+\delta$. For $b\le B$ each valuation in this sum may be capped at $\Theta$
without changing the answer: any summand reaching $\Theta$ already suffices, since all other
summands are nonnegative. By \eqref{cost:eq:cap} every lift has the weight of its parent
divided by $Q$. On $\Sigma'$ use the cap $\kappa=\Theta/Q$. Then \eqref{cost:eq:capdiff}
identifies every capped comparison on $\Sigma'$ with the corresponding comparison on
$\Sigma$, divided by $Q$. Sending marks to their parents preserves multiplicities;
conversely one may choose any lift of each occurrence of a mark, and testing a target at
all its lifts repeats the test at its parent. Both the optimum and the degree comparison
follow.
\end{proof}

Iteration divides cost and tolerance by $Q$ at each step.

\begin{example}[An elliptic curve with overlapping torsion valuations]\label{cost:ex:elliptic}
Let $K=\Q_2(\pi)$ with $\pi^4=2$, and consider
\[
 E\colon\quad y^2+\pi xy+y=x^3 .
\]
Its discriminant $\pi^3-27$ is a unit. On the special fibre negation sends $(x,y)$ to
$(x,y+1)$, so the only geometric point killed by $2$ is the identity; the reduction is
supersingular, and all geometric $2$-power torsion lies in the formal group, which has
height two, so $Q=4$. With $t=-x/y$ and $u=-1/y$ the equation becomes
$u(1-\pi t-u)=t^3$, so $u=t^3+O(t^4)$ in $\OO_K[[t]]$, and the invariant differential is
\[
 \frac{\dd t}{1-\pi t-2u}=\bigl(1+\pi t+\pi^2t^2+(\pi^3+2)t^3+O(t^4)\bigr)\,\dd t.
\]
Comparing coefficients through degree four in $L([2](t))=2L(t)$ for the formal logarithm
$L$ gives
\[
 [2](t)=2t-\pi t^2-7t^4+O(t^5),
\]
with all coefficients integral; see~\cite[Ch.~IV, \S\S1--5]{Silverman2009Arithmetic} for
these constructions. Here $\kappa=1/8$ and $\Theta=1/2$.

\emph{Valuations at every level.} The negative-slope part of the Newton polygon of
$[2](T)$ has vertices $(1,1)$, $(2,1/4)$, $(4,0)$: the three nonzero points killed by $2$
have valuations $3/4$ once and $1/8$ twice. If $c$ is the coordinate of the point of
valuation $3/4$, the polygon of $[2](T)-c$ has vertices $(0,3/4)$, $(2,1/4)$, $(4,0)$, so
two of its preimages have valuation $1/4$ and two have valuation $1/8$. For a nonzero target
of valuation $0<y\le1/2$ the relevant part of the polygon is a single segment from $(0,y)$
to $(4,0)$, and all four preimages have valuation $y/4$. Hence, for every $n\ge2$, with
$s=4^{2-n}$, the points of exact order $2^n$ have the valuations
\[
\begin{array}{c|ccc}
v(t) & s/4 & s/8 & s/32 \\
\hline
\text{number of points} & 2\cdot4^{n-2}&2\cdot4^{n-2}&8\cdot4^{n-2}.
\end{array}
\]
Levels one and two share the valuation $1/8$, and for every $n\ge2$ the least valuation at
level $n$ equals the middle valuation at level $n+1$.

\emph{The tree at exact order four.} Let $\Gamma=E[4]$, $\Gamma'=E[2]$,
$\Sigma=\Gamma\setminus\Gamma'$ and $\Gamma_r=\set{P\in\Gamma:v(t(P))\ge r}$. The counts are
\[
\begin{array}{c|rrrrr}
r & 1/32 & 1/8 & 1/4 & 3/4 & +\infty\\
\hline
\abs{\Gamma_r}       &16&8&4&2&1\\
\abs{\Gamma'\cap\Gamma_r} &4&4&2&2&1 .
\end{array}
\]
The two cosets of $\Gamma_{1/8}$ contain four and eight points of $\Sigma$. In the coset of
the origin each of the two cosets of $\Gamma_{1/4}$ contains two points of $\Sigma$, forming
one coset of $\Gamma_{3/4}$. In the other coset of $\Gamma_{1/8}$ all eight points survive;
its two cosets of $\Gamma_{1/4}$ each split into two pairs modulo $\Gamma_{3/4}$, and each
pair splits into two leaves.

\emph{Five roots suffice at every exact order from four on.} Fix $n\ge2$, test all points of
exact order $2^n$ at tolerance zero with $\rho\ge s/4$. Then the quantities of
Theorem~\ref{cost:thm:finite} are
\[
 d_*=5,\qquad B_*=\frac{15s}{32}=\frac{15}{32\cdot4^{n-2}}.
\]
In particular at order four a budget strictly below $15/32$ ensures a good point, whereas a
quintic of budget exactly $15/32$ blocks every point. To see this, choose a compatible
torsion basis $P_n,Q_n$: at level one take the point of valuation $3/4$ and either point of
valuation $1/8$, lift the first to valuation $1/4$ at level two, and continue by arbitrary
preimages. Odd multiples preserve valuation. On the two cyclic directions the valuations at
exact order $2^j$ are
\[
 a_1=3/4,\quad a_j=4^{1-j}\ (j\ge2),\qquad b_j=\frac1{8\cdot4^{j-1}}\ (j\ge1).
\]
The two lists are disjoint, so the valuation of $aP_n+bQ_n$ is the smaller of the two
component valuations, and this computes every pairwise distance. The two inner parts
consist of the points with $a$ odd and $b\equiv0$, respectively $b\equiv2\pmod4$. Within
either part, differences have $a$ even and $b$ divisible by four, hence valuation at least
$s/2$. One mark therefore exceeds the threshold $0$ on the part of weight $s/4$ and the
threshold $s/8$ on the part of weight $s/8$. The outer part, $b$ odd, has weight $s/32$. It
splits by $b\bmod4$ into two branches at mutual distance $s/8$, and each splits by
$a\bmod2$ with cross valuation $s/4$. Its greedy sequence has $E(1)=3s/32$ and
$E(2)=10s/32$, and a third mark adds at least $3s/32$ wherever it matters. The threshold
after the first two parts is exactly $10s/32$, so the strict inequality costs three outer
marks. The total is $s/4+s/8+3s/32$. For $n>2$ the same values follow from $n=2$ by
Proposition~\ref{cost:prop:scaling}, since $15/32<\Theta$.

At exact order four and $\rho=1$ the recursion \eqref{cost:eq:rec} gives the following
values at positive tolerances.
\begin{center}
\begin{tabular}{@{}crr@{}}\toprule
Tolerance & $d_*$ & $B_*$\\
\midrule
$0$ &5&$15/32$\\
$1/32$ &5&$15/32$\\
$1/8$ &6&$1/2$\\
$1/2$ &9&$29/32$\\
$1$ &12&$1$\\
\bottomrule\end{tabular}
\end{center}

\emph{Two overlapping levels.} For the union of the exact orders $8$ and $16$, in
coordinates $(a,b)\in(\Z/16)^2$ with respect to a basis chosen as above, the five parts in
order of decreasing weight are described by
\begin{gather*}
 (a\equiv2\bmod4,\ b\equiv0\bmod8),\qquad
 (a\equiv2\bmod4,\ b\equiv4\bmod8),\\
 (a\text{ odd},\ 4\mid b),\qquad b\equiv2\bmod4,\qquad b\text{ odd},
\end{gather*}
and the recursion at tolerance zero runs as follows.
\[
\begin{array}{c|r|c|r|cc}
\omega& \abs{\Sigma_\omega}& t_\omega&m_\omega&E_\omega(m_\omega-1)&E_\omega(m_\omega)\\\hline
1/16 &8&0&1&0&1/16\\
1/32 &8&1/32&1&0&3/32\\
1/64 &32&1/16&3&1/16&5/64\\
1/128 &64&13/128&6&13/128&1/8\\
1/512 &128&85/512&12&85/512&25/128
\end{array}
\]
Hence $B_*=27/128$, attained with $23$ roots. By Proposition~\ref{cost:prop:scaling}, for every
$n\ge4$ the union of the exact orders $2^{n-1}$ and $2^n$ has $B_*=27/(128\cdot4^{n-4})$
and $d_*=23$. The exact-arithmetic outputs are in the following files; the first also
compares the greedy values with an exhaustive enumeration of multisets of marks on the
order-four test cases.
\costfiles{
\item \artifact{grh_method_transfers_research/code/torsion_checks.txt}
\item \artifact{grh_method_transfers_research/code/two_level_checks.txt}}
\end{example}

\begin{proposition}[Eventual scaling for every fixed pattern; \stU]\label{cost:prop:eventual}
Fix $\rho>0$ and a finite nonempty pattern of $k$ distinct positive exact-order exponents.
Translate all exponents upward by the same integer and take tolerance zero. From an explicit
level onward, each unit translation divides the least blocking budget by $Q$ and preserves
the least degree of a budget-minimizing blocker. With
\[
 \Theta_*=\min\set{\Theta,\rho},\qquad
 L=1+\floor*{\frac1{\min\set{1,\Theta_*}}},
\]
it suffices that the smallest exponent $j_{\min}$ of the translated pattern satisfies
$j_{\min}\ge L+k+2+s$ for some integer $s\ge0$ with $Q^s>k/\Theta_*$.
\end{proposition}

\begin{proof}
Write $\Gamma_j=F[p^j]$ and $Y_j=\Gamma_j\setminus\Gamma_{j-1}$. Subtracting successive
levels in \eqref{cost:eq:totalval} gives $\sum_{y\in Y_j}v(y)=1$. For $x\in Y_j$ the product
of the differences against all $y\in\Gamma_j$ other than $x$ again has valuation $j$, by
translation invariance of the distance. Evaluating the prepared factorization of
$[p^{j-1}]$ at $x$ gives the row identity
\[
 \sum_{y\in Y_j,\ y\ne x}v(x-y)=j-v([p^{j-1}]x)\ge j-1,
\]
because a nonzero point of $\Gamma_1$ has valuation at most the total of its shell, which
is one.

Every torsion point of valuation at least $\Theta_*$ belongs to $\Gamma_L$. Indeed for a
point of valuation $v_0\ge\Theta_*$ a nonzero multiple by $p$ has valuation at least
$v_0+\min\set{1,\Theta_*}$, by the integral multiplication series and its linear term $pT$;
iterating to the last nonzero multiple, which has valuation at most one, gives the stated
bound on the order. Therefore capping all off-diagonal pair valuations in any one row at
$\Theta_*$ costs at most $L$ in total: the formal differences are distinct nonzero torsion
points, and those whose valuations exceed $\Theta_*$ belong to $\Gamma_L$, where the sum of
all nonzero valuations is $L$.

Choose a parent pattern whose smallest exponent is at least $L+k+2$. All its points have
valuation strictly below $\Theta_*$. Mark every parent point once; the cost is $k$. At any
parent target, its own exact-order shell alone contributes an off-diagonal capped row sum of
at least $j-1-L>k$. Now lift the pattern by $s$ exponents and choose one arbitrary lift
under $[p^s]$ of each marked parent. Iterating \eqref{cost:eq:capdiff}, which also holds
with $\Theta$ replaced by any smaller positive cap, gives for all points $x,y$
\[
 \min\set{v([p^s]x-[p^s]y),\Theta_*}=Q^s\min\set{v(x-y),\Theta_*/Q^s}.
\]
Each chosen mark has the weight of its parent divided by $Q^s$, and every target of the
lifted pattern lies over a parent target, so its capped row sum exceeds $k/Q^s$. This is a
blocker of cost $k/Q^s<\Theta_*$. The optimal cost is then also below $\Theta_*$, and
Proposition~\ref{cost:prop:scaling} applies at each further level. All these target and mark valuations
stay below $\rho$, as required for the fixed radius.
\end{proof}

\subsection{Nonvanishing alone}\label{tors:sec:nonvanishing}
\sources{grh_beyond_dirichlet_grh_2026-11-25; grh_quantitative_twists_2026-11-25; dirichlet_li_recovery}

When an application needs only $F_i(z)\ne0$, the interval of exponents can be shorter
than in \cref{tors:thm:selection}: its length is $\floor{B}$ instead of
$\floor{Bp^{h_0-1}}$. The proof is different. It divides by the polynomial whose roots
are the points of the grid, and it treats all variables at once. The same method, applied
to a product over the unit group, determines exactly the budget needed to force a zero at
one of a prescribed set of unit powers on a shell (\cref{tors:thm:cover}).

\begin{lemma}[Integral division by a distinguished polynomial]\label{tors:lem:division}
Let $g(T)=T^d+a(T)\in\Zp[T]$ be monic of degree $d\ge1$, with $\deg a<d$ and every
coefficient of $a$ divisible by $p$. Every $F\in\OO[[T_1,\dots,T_s]]$ has a decomposition
\[
  F=g(T_1)\,Q+R,\qquad Q,R\in\OO[[T_1,\dots,T_s]],\qquad \deg_{T_1}R<d .
\]
Successive division in each variable gives
\begin{equation}\label{tors:eq:multidivision}
  F=\sum_{a=1}^sg(T_a)\,Q_a+R,\qquad \deg_{T_a}R<d\quad(1\le a\le s),
\end{equation}
with all $Q_a$ integral series and $R$ an integral polynomial.
\end{lemma}

\begin{proof}
Bounded coefficient arrays with the supremum norm of the coefficients form a complete
normed space over $\Cp$, whose unit ball is the set of integral series. Let $S_d$ discard
the terms of degree less than $d$ in $T_1$ and lower the remaining powers of $T_1$ by
$d$, and let $M_a$ be multiplication by $a(T_1)$. The operator $K=S_dM_a$ has norm at
most $p^{-1}$, because every coefficient of $a$ has valuation at least one. Therefore
\[
  Q=\sum_{n=0}^\infty(-K)^nS_dF
\]
converges to an integral series and satisfies $(1+K)Q=S_dF$. For $R=F-g(T_1)Q$ this
identity says $S_dR=0$, that is, $\deg_{T_1}R<d$. Division of $R$ by $g(T_2)$ preserves
the degree bound in $T_1$, since the operators then act on $T_2$ only. Continuing through
all variables gives \eqref{tors:eq:multidivision}.
\end{proof}

An integral series $P$ \emph{vanishes to order at least $M$} at a point $a$ of the open
unit polydisc if all terms of total degree less than $M$ of its expansion in the increments
$T_b-a_b$ vanish. A product of $M$ integral series that vanish at $a$ vanishes to order at
least $M$ at $a$. The following multiplicity form of the grid lemma is the analogue for
bounded series of the identity of Ball and Serra for polynomials over a field
\cite[Theorem~3.1]{BallSerra2009Punctured}; here integral quotients are needed to retain
the Gauss bound.

\begin{lemma}[Vanishing on a finite product grid]\label{tors:lem:grid}
In \cref{tors:lem:division} suppose that $g$ has $d$ distinct roots, forming a set
$A\subset\Cp$, all of positive valuation, and let $M\ge1$. If an integral series $P$
vanishes to order at least $M$ at every point of $A^s$, then
\[
  P\in\bigl(g(T_1),\dots,g(T_s)\bigr)^M\,\OO[[T_1,\dots,T_s]],\qquad
  G_\rho(P)\ge M\,G_\rho(g)\quad(\rho>0).
\]
For $M=1$: an integral series vanishing on $A^s$ lies in $(g(T_1),\dots,g(T_s))$.
\end{lemma}

\begin{proof}
The polynomial $g^M$ satisfies the hypotheses of \cref{tors:lem:division}, since
$g^M\equiv T^{dM}\pmod p$. Dividing successively in each variable by $g(T_b)^M$ gives
$P=\sum_bg(T_b)^MQ_b+R$ with integral quotients and an integral polynomial $R$ of degree
less than $Md$ in each variable. Expand $R$ uniquely by repeated monic division as
\[
  R=\sum_{0\le\beta_b<M}S_\beta(T)\prod_bg(T_b)^{\beta_b},\qquad \deg_{T_b}S_\beta<d ,
\]
with integral $S_\beta$. The terms $g(T_b)^MQ_b$ vanish to order at least $M$ on $A^s$, so
$R$ does too. Induct on $k<M$ to show $S_\beta=0$ for $\abs\beta=k$. At $a\in A^s$, after all
terms with $\abs\beta<k$ have disappeared, the homogeneous term of degree $k$ in the
increments $X_b=T_b-a_b$ is $\sum_{\abs\beta=k}S_\beta(a)\prod_b(g'(a_b)X_b)^{\beta_b}$.
Distinct monomials and simple roots force $S_\beta(a)=0$ for all such $\beta$, and
interpolation in one coordinate at a time gives $S_\beta=0$ identically. The surviving
terms have $\abs\beta\ge M$ and belong to the asserted power of the ideal. All multipliers
are integral, so \cref{tors:lem:gauss}(a) gives the Gauss bound.
\end{proof}

\paragraph{Unit powers and their covering densities.}
For $h\ge1$ put $\torsU_h=(\Z/p^h\Z)^\times$. Let $S$ be a nonempty finite set of levels, and put
$A_S=\bigcup_{h\in S}Z_h$, $H=\max S$ and $N=\#S$. Fix $s\ge1$, a depth
$r\ge\max_{h\in S}r_h$, and a finite nonempty set $J$ of positive integers prime to $p$.
For $z$ in the open unit polydisc and an integer $j$ write
$[j]z=\bigl((1+z_b)^j-1\bigr)_{b=1}^s$; thus $F([j]z)=F^{[j]}(z)$ in the notation of
\cref{tors:lem:powers}. An \emph{admissible family} is a finite collection of
$F_i\in\OO[[T_1,\dots,T_s]]$ with witnesses $x_i\in\Cp^s$ of depth $r$ and finite total
budget $B=\sum_iv(F_i(x_i))$. It \emph{blocks $A_S^s$ at the powers $J$} if for every
$z\in A_S^s$ there are $i$ and $j\in J$ with $F_i([j]z)=0$. Let $C_s(S,J)$ be the infimum
of the budgets of blocking families, with $r$ fixed. Let $J_h$ be the image of $J$ in
$\torsU_h$ and define the covering number and the covering density
\[
  \tau_h=\min\bigl\{\#D:\ D\subseteq\torsU_h,\ (uJ_h)\cap D\ne\varnothing\text{ for every }
  u\in\torsU_h\bigr\},\qquad c_h=\frac{\tau_h}{e_h}.
\]
Thus $DJ_h^{-1}=\torsU_h$. This is an ordinary finite covering number; exact factorizations
of finite groups and the related $k$-radius sequence problem have their own literature
\cite{BlackburnMcKee2012Constructing}.

\begin{theorem}[Exact cost of blocking nonvanishing at unit powers; \stU]\label{tors:thm:cover}
In the setting above, for a single shell $C_s(\set{h},J)=c_h$ in every dimension $s$. For
every $S$,
\[
  C_1(S,J)=\sum_{h\in S}c_h,\qquad
  \max_{t\in S}\bigl(\#(S\cap[1,t])\,c_t\bigr)\ \le\ C_s(S,J)\ \le\ \sum_{h\in S}c_h .
\]
In particular $C_s(S,J)=Nc$ in every dimension if $c_h=c$ for all $h\in S$. If $B<C_s(S,J)$,
every admissible family of budget at most $B$ has a tuple $z\in A_S^s$ with
$F_i([j]z)\ne0$ for all $i$ and all $j\in J$. In every exact formula above a blocking
family exists at equality; it can be taken to be one polynomial in the first coordinate
with the witness $0$. All formulas hold for every permitted depth $r$.
\end{theorem}

\begin{proof}
Substitution by $[u]$ for $p\nmid u$ preserves $G_r$: the series $(1+T)^u-1$ and its
compositional inverse $(1+T)^{1/u}-1$ have integral coefficients and zero constant term,
because $1/u\in\Zp$ and binomial coefficients are continuous and integral on $\Zp$, so
neither substitution can decrease $G_r$.

Let $F=\prod_iF_i$ for a blocking family. The group $\torsU_H$ acts diagonally on $A_S^s$ by
$z\mapsto[u]z$. The orbit of a tuple whose largest coordinate order exponent is $t$ is a copy
of $\torsU_t$: a stabilizing element must be $1\bmod p^t$ because of the coordinate of that
exact order, and this condition fixes the whole tuple. On the orbit the positions where $F$
vanishes meet every translate $uJ_t$, so they number at least $\tau_t$. The inverse image
of a cover under the reduction $\torsU_H\to\torsU_t$ is again a cover, hence $c_H\le c_t$.

Choose a positive integer representative $u$ for each class in $\torsU_H$ and put
$P(T)=\prod_{u\in\torsU_H}F([u]T)$. At each tuple of $A_S^s$ at least
$(e_H/e_t)\tau_t\ge\tau_H$ factors vanish, so $P$ vanishes there to order at least
$\tau_H$. Apply \cref{tors:lem:grid} with $g=\prod_{h\in S}\Phi_{p^h}(1+T)$, which is monic,
congruent to a power of $T$ modulo $p$, and has the distinct roots $A_S$; by
\eqref{tors:eq:gauss} $G_r(g)=\sum_{h\in S}e_hr_h=N$. With \cref{tors:lem:gauss} and
\eqref{tors:eq:witness},
\[
  e_HB\ \ge\ e_HG_r(F)=G_r(P)\ \ge\ N\tau_H ,
\]
that is $B\ge Nc_H$. The same argument on $S\cap[1,t]$ proves each prefix bound.

Conversely, choose a minimal cover $D_h\subseteq\torsU_h$ for each $h\in S$ and a root of unity
$\xi_h$ of order $p^h$, and let
\[
  F_*(T)=\prod_{h\in S}\prod_{u\in D_h}\bigl(T_1-(\xi_h^u-1)\bigr).
\]
The first coordinate of every tuple of $A_S^s$ has a power $[j]$, $j\in J$, in the root
set. At the witness $0$ the valuation is $\sum_h\#D_h\,r_h=\sum_hc_h$. In one variable a
blocking product must itself vanish at $\tau_h$ or more distinct points of each shell
$Z_h$, $h\in S$. An integral series vanishing at a point $a$ of positive valuation is
divisible by $T-a$ with integral quotient, as follows from summing $(T^n-a^n)/(T-a)$
termwise, and distinct roots persist after each division. Every extracted factor has
$G_r(T-a)=r_h$, which proves the one-variable lower bound $\sum_hc_h$. All the stated
equalities follow.
\end{proof}

\begin{corollary}[Explicit simultaneous nonvanishing; \stU]\label{tors:thm:nonvanishing}
Let $s\ge1$, and consider integral series $F_1,\dots,F_n$ in the variables
$T_1,\dots,T_s$. Let $x_i\in\Cp^s$ be witnesses of depth $r>0$ for the $F_i$ such that
\[
  \sum_{i=1}^nv(F_i(x_i))\le B<\infty .
\]
Put $H=h_*(r)+\floor{B}$. There are roots of unity $\zeta_1,\dots,\zeta_s$, with
$\zeta_a$ of exact order $p^{h_a}$ and $h_*(r)\le h_a\le H$, such that
\[
  F_i(\zeta_1-1,\dots,\zeta_s-1)\ne0\qquad(1\le i\le n).
\]
Thus $\zeta_a^{p^H}=1$ and $0<v(\zeta_a-1)\le r$. The exponent $H$ depends only on
$p,r,B$; it does not depend on the number of variables, on the number of functions
separately from their total budget, or on degrees. It is optimal in every dimension.
\end{corollary}

\begin{proof}
Take $J=\set1$ in \cref{tors:thm:cover}. Then $\tau_h=e_h$ and $c_h=1$ for every $h$, so
the window $S=[h_*(r),H]$ of $N=\floor{B}+1>B$ shells has $C_s(S,\set1)=N$. If
$\floor B\ge1$, the window of the first $\floor B$ of these shells is blocked at budget
$\floor B\le B$, and if $\floor B=0$, the level $h_*(r)$ is the least one allowed.
\end{proof}

The case $M=1$ of \cref{tors:lem:grid} already gives the corollary directly: a product $F$
vanishing on $\torsZ{h_*}{H}^s$ would satisfy
$G_r(F)\ge G_r\bigl(\prod_{h=h_*}^H\Phi_{p^h}(1+T)\bigr)=H-h_*+1>B$.

\begin{remark}[Sharpness and comparison]\label{tors:rem:nonvanishing}
(1) The blocker of \cref{tors:thm:cover} for $J=\set1$ is explicit. Put $k=\floor{B}$. If
$k\ge1$, the polynomial $F(T)=\prod_{h=h_*}^{h_*+k-1}\Phi_{p^h}(1+T)$ has the witness $x=0$
with $v(F(0))=v(p^k)=k\le B$, and it vanishes at every point of $\torsZ{h_*}{h_*+k-1}$. If
$k=0$, the level $h_*$ is the least one allowed by the radial condition. More generally,
for $h_0\le L$ and $\rho\ge r_{h_0}$ the polynomial $\prod_{h=h_0}^{L}\Phi_{p^h}(1+T)$ has
$G_\rho=L-h_0+1$ and vanishes on $\torsZ{h_0}{L}$, so a linear dependence of the exponent
on the budget is unavoidable. For the version with several powers see
\cref{tors:thm:twopowers}.

(2) For $h_*(r)=1$, that is for $r\ge1/(p-1)$, \cref{tors:thm:nonvanishing} is contained
in \cref{tors:cor:packets}(a), which gives the same exponent $1+\floor{B}$ together with
the bound $\sum_iv(F_i(z))\le B$. For $h_*(r)\ge2$ the two results are different: forcing
$F$ to \emph{vanish} on a shell costs a budget $1$, whereas forcing
$v(F(z))>G_\rho(F)$ on a shell costs only $p^{1-h_0}$ by \cref{cost:thm:minimax}\textup{(i)}
with $H=h_0$.

(3) \Cref{tors:thm:nonvanishing} guarantees finite values, not values close to the
witness valuations; that stronger conclusion is the subject of \cref{tors:sec:tolerance}.

(4) In several variables a tuple may have coordinates of different orders, and its unit
orbit is controlled by the largest of them; this is why the lower bound of
\cref{tors:thm:cover} for mixed shells uses $c_t$ at a prefix endpoint. Powers divisible by
$p$ move points between shells and lead to a different covering problem.
\end{remark}

\subsection{Packets with individual thresholds over unrestricted coefficients}\label{cost:sec:packet}
{\raggedright\sources{grh-torsion-and-prime-sampling; grh-torsion-and-zero-data; grh_exact_torsion_transfer_2026-11-25; grh_sharp_transfers_2026-11-25}}

We return to setting (C) and to the packet convention of Definition~\ref{cost:def:packet}. Here every
shell is charged separately, and the charge of a shell depends only on the capacity
function \eqref{cost:eq:capacity}.

\begin{lemma}[Charging a shell]\label{cost:lem:charge}
Let $(g_i,t_i)_i$ be a packet over $\OO$ and let $S$ be a finite set of positive integers
with $\rho\ge\max_{h\in S}r_h$. For $h\in S$ let $M_h$ be the number of roots of valuation
$r_h$ of all the $g_i$ together, with multiplicity. Then
\begin{equation}\label{cost:eq:charge}
 \mathcal L(z)\le\sum_{v(\alpha)=r_h}\bigl(v(z-\alpha)-r_h\bigr)\quad(z\in Z_h),\qquad
 \sum_{h\in S}r_hM_h\le\sum_i\Delta_\rho(g_i)\le\sum_it_i,
\end{equation}
where the first sum runs over those roots.
\end{lemma}

\begin{proof}
Since $t_i\ge G_\rho(g_i)\ge G_{r_h}(g_i)$, each positive part is at most
$v(g_i(z))-G_{r_h}(g_i)$, which is given by \eqref{cost:eq:shellvalue}. The second chain is
\eqref{cost:eq:budget} together with $G_0(g_i)\ge0$.
\end{proof}

For $\delta\ge0$ and $h\ge1$ put
\begin{equation}\label{cost:eq:shellcost}
 m_h(\delta)=\min\set{0\le m\le e_h:U_h(m)>\delta},\qquad c_h(\delta)=\frac{m_h(\delta)}{e_h}.
\end{equation}
The strict inequality is essential when $\delta$ is an attained loss.

\begin{theorem}[Exact obstruction budget for packets; \stU]\label{cost:thm:packet}
Let $h_0\ge1$, $H\ge h_0$, $\rho\ge r_{h_0}$ and $\delta\ge0$. The least cost of a packet
over $\OO$ blocking every shell from $h_0$ through $H$ at tolerance $\delta$ is
\[
 C_\OO(\set{h_0,\ldots,H},\delta)=\sum_{h=h_0}^Hc_h(\delta).
\]
The minimum is attained by monic linear polynomials with common witness $0$ and
coefficients in the ring of integers of $\Qp(\zeta_{p^H})$. At most $H-h_0+1$ polynomials
suffice if their degrees are unrestricted.
\end{theorem}

\begin{proof}
Let a packet block the shell $h$. By \eqref{cost:eq:charge} the loss at $z\in Z_h$ is at
most the on-shell excess of the $M_h$ roots of valuation $r_h$. If $M_h<m_h(\delta)$, then
$M_h<e_h$ and Lemma~\ref{cost:lem:shell}(i) gives a point with loss at most
$U_h(M_h)\le\delta$. Hence $M_h\ge m_h(\delta)$ on every blocked shell, and the second part
of \eqref{cost:eq:charge} gives the lower bound.

For equality, on each shell mark the first $m_h(\delta)$ positive integers $b$ prime to
$p$ and insert the polynomial $T-(\xi_h^b-1)$ with threshold $r_h$, which is both its value
at $0$ and its Gauss valuation at $\rho$. By Lemma~\ref{cost:lem:shell}(iii) the loss at an
unmarked point of the shell is at least $U_h(m_h(\delta))>\delta$, and a marked point gives
an infinite loss; this includes the case $m_h(\delta)=e_h$. Linear factors from other
shells have zero positive excess. The cost is the claimed sum. Finally, on each shell one
may multiply its linear factors into one polynomial, since on that shell each individual
excess is nonnegative and on every other shell the excesses are nonpositive.
\end{proof}

\begin{remark}[Moving baseline]\label{cost:rem:moving}
For $h_0=1$ and $\rho=1$ the same number $\sum_{h=1}^Hc_h(\delta)$ is the least value of
$G_1(g)$ of a single integral polynomial $g$ with
$v(g(z))>G_{r_h}(g)+\delta$ for every $z\in Z_h$ and every $1\le h\le H$, that is, when the
excess on each shell is measured over the moving Gauss valuation $G_{r_h}$. The lower bound
is the same argument, because the excess over $G_{r_h}$ is exactly the on-shell excess
\eqref{cost:eq:shellvalue}. The product of the extremal packet attains it, because its
excess is additive and each summand is nonnegative.
\end{remark}

At tolerance zero the costs are explicit. For $h\ge2$ the capacity $U_h(m)$ first becomes
positive at $m=p-1$, where $U_h(p-1)=p^{1-h}$, and for $h=1$ the whole shell must be
marked. Thus $c_h(0)=p^{1-h}$ for every $h\ge1$ and
\begin{equation}\label{cost:eq:zerotol}
 C_\OO(\set{h_0,\ldots,H},0)=\sum_{h=h_0}^Hp^{1-h}<B_0:=\frac{p^{2-h_0}}{p-1}.
\end{equation}

\begin{corollary}[Individual and moving comparisons have no endpoint; \stU]\label{cost:cor:individual}
Let $h_0\ge1$ and $\rho\ge r_{h_0}$. For every $H\ge h_0$ there is a finite packet of monic
linear polynomials $g_i$ with $\sum_iG_\rho(g_i)<B_0$ such that at every point $z$ of the
shells $h_0,\ldots,H$ some member satisfies $v(g_i(z))>G_\rho(g_i)$. For $p=2$ one may take the polynomials $T-(\xi_h-1)$,
$h_0\le h\le H$, of total budget $\sum_hr_h<2r_{h_0}$; their product $g$ satisfies, for
$h_0=1$ and $\rho=1$, $G_1(g)=2-2^{1-H}<2$ and $v(g(z))>G_{r_h}(g)$ on every shell
$1\le h\le H$.
\end{corollary}

\begin{proof}
The first statement is \eqref{cost:eq:zerotol}. For $p=2$ one has $c_h(0)=r_h$ and one mark
per shell; two primitive roots of the same $2$-power order have a quotient of strictly
smaller order unless they are equal. The last statement is Remark~\ref{cost:rem:moving}.
\end{proof}

Thus a bounded budget can defeat every individual exact comparison, and every comparison
with the moving target $G_{r_h}$, on arbitrarily many shells, whereas by
Corollary~\ref{cost:cor:endpoint} the comparison of the product with its fixed Gauss valuation
$G_\rho$ succeeds by the shell $h_0+\floor{Bp^{h_0-1}}$: the savings on the earlier factors
offset the loss on the active one. This is the reason why \cref{tors:thm:selection} is
stated for an aggregate valuation and a fixed radius.

At positive tolerance the cost of a shell decays like the harmonic series.

\begin{corollary}[Exponential depth is necessary; \stU]\label{cost:cor:depth}
For fixed $p$, $h_0$ and $\delta>0$,
\[
 c_h(\delta)=\frac\delta h+O_{p,\delta}\Bigl(\frac{\log h}{h^2}\Bigr),\qquad
 \sum_{h=h_0}^Hc_h(\delta)=\delta\log H+O_{p,h_0,\delta}(1).
\]
The least $H=H_*(B,\delta)$ such that every packet over $\OO$ of cost at most $B$ has a point
of loss at most $\delta$ on some shell $h_0\le h\le H$ is the first $H$ with
$\sum_{h=h_0}^Hc_h(\delta)>B$. Consequently
$\log H_*(B,\delta)=B/\delta+O_{p,h_0,\delta}(1)$ as $B\to\infty$.
\end{corollary}

\begin{proof}
For $m<e_h$ one has $U_h(m)\le(h-1)m/e_h$, so $c_h(\delta)\ge\delta/(h-1)$ for large $h$,
also when $m_h(\delta)=e_h$. For the reverse bound let
$\ell=h-\ceil{3\log_ph}$; eventually $1\le\ell<h$. For any $m<e_h$,
\[
 U_h(m)\ge\frac{\ell m-\sum_{j=1}^\ell e_j}{e_h}=\frac{\ell m-p^\ell+1}{e_h}.
\]
Take $m=\floor{(\delta e_h+p^\ell)/\ell}+1$, which is eventually smaller than $e_h$; then
$U_h(m)>\delta$. Since $p^\ell/e_h=O_p(h^{-3})$ and $h-\ell=O_p(\log h)$, this gives
$c_h(\delta)\le\delta/h+O_{p,\delta}(\log(h)/h^2)$. Summation and the convergence of the
error series give the formula for the sum. The strict inequality in the criterion for
$H_*$ holds because the minimum in Theorem~\ref{cost:thm:packet} is attained.
\end{proof}

In particular no bound on the order exponent that is polynomial in the budget and in the
inverse tolerance, uniformly in the number and the degrees of the polynomials, can hold for
the total positive excess: the order exponent itself must be of size $\exp(B/\delta+O(1))$.

The dependence on the tolerance at a fixed budget has a threshold at $B_0$.

\begin{theorem}[The fixed-budget transition; \stU]\label{cost:thm:transition}
Let $s_h=p^{1-h}$, so that the number $B_0$ of \eqref{cost:eq:zerotol} is
$\sum_{h\ge h_0}s_h=p^{2-h_0}/(p-1)$. For fixed $p$, $h_0$ and
$B\ge0$, the least guaranteed exponent $H_*(B,\delta)$ of Corollary~\ref{cost:cor:depth} obeys the
following laws as $\delta\downarrow0$.
\begin{enumerate}[label=\textup{(\roman*)}]
\item If $B<B_0$, it eventually equals the first $H$ with $\sum_{h=h_0}^Hs_h>B$.
\item If $B=B_0$, put $t=\ceil{1+\log_p(1/\delta)}$. For all sufficiently small
$\delta>0$, $H_*(B_0,\delta)=t$ for odd $p$ and $H_*(B_0,\delta)=t+1$ for $p=2$.
\item If $B>B_0$, then $\lim_{\delta\downarrow0}\delta\log H_*(B,\delta)=B-B_0$.
\end{enumerate}
\end{theorem}

\begin{proof}
For each fixed shell, $c_h(\delta)=s_h$ for every sufficiently small $\delta\ge0$, by the
computation before \eqref{cost:eq:zerotol}. This proves (i) and shows that
$H_*(B,\delta)\to\infty$ when $B\ge B_0$.

The estimates in the proof of Corollary~\ref{cost:cor:depth} give, for $h$ sufficiently large and
$0\le\delta\le1$, the uniform summable bound
\[
 0\le c_h(\delta)-\frac\delta h\ll_p\frac{\log(2h)}{h^2};
\]
the upper bound uses $c_h\le\delta/\ell+p^\ell/(\ell e_h)+1/e_h$ and the lower bound uses
$U_h(m)\le(h-1)m/e_h$. Dominated convergence on the positive integers yields, along any
$H\to\infty$, $\delta\downarrow0$,
\[
 \sum_{h=h_0}^Hc_h(\delta)-\delta\log H\longrightarrow B_0,
\]
since $\sum_{h\le H}1/h-\log H$ is bounded. At $H=H_*$ the overshoot above $B$ is at most
$c_H(\delta)$ once $H>h_0$, and this tends to zero. This proves (iii).

For (ii) take $\delta$ so small that $t>\max\set{2,h_0}$. All earlier costs equal $s_h$. At
$h=t$ the tolerance lies in $[s_t,ps_t)$. All $e_j$ are multiples of $p-1$, so the capacity
is constant between consecutive multiples of $p-1$ in the number of marks; hence
$c_t(\delta)\ge2s_t$. The remaining baseline tail is $\sum_{h>t}s_h=s_t/(p-1)$. For $p>2$
the extra cost already exceeds this tail and $H_*=t$. For $p=2$ the capacities at $m=1,2$
are $s_t$ and $4s_t$, since $t\ge3$; thus $c_t(\delta)=2s_t$, the cumulative cost is exactly
$B_0$, and one further shell of positive cost makes it strictly larger.
\end{proof}

For $p=2$ and $h_0=1$ the critical budget is $2$. Raising the budget from exactly $2$ to
$3$ changes the least guaranteed exponent from $\ceil{1+\log_2(1/\delta)}+1$ to
$\exp((1+o(1))/\delta)$. So the failure of a bound polynomial in $1/\delta$ and independent
of the packet size already occurs at a fixed budget.

\begin{example}\label{cost:ex:packet}
At $p=2$, $h_0=1$ and $\delta=1/4$ the optimal costs are as follows.
\begin{center}\begin{tabular}{@{}lrrrrrr@{}}\toprule
Shell exponent $h$ & 1 & 2 & 3 & 4 & 5 & 6\\
\midrule
Primitive roots $e_h$ & 1&2&4&8&16&32\\
Minimum marks $m_h(1/4)$ &1&1&2&2&3&4\\
Cumulative budget & $1$ & $3/2$ & $2$ & $9/4$ & $39/16$ & $41/16$\\
\bottomrule\end{tabular}\end{center}
A budget of exactly $2$ can block the orders $2$, $4$, $8$, and every packet of that budget
has a point of loss at most $1/4$ of order at most $16$. At order $8$ one mark attains
minimum loss exactly $1/4$; equality counts as success, so blocking this shell requires a
second mark. At $p=2$, $h_0=1$, $\delta=1/8$ and $B=3$ the exact cut-off is
\[
 2.9998<\sum_{h=1}^{371}c_h(1/8)<3<\sum_{h=1}^{372}c_h(1/8)<3.0003,\qquad H_*(3,1/8)=372:
\]
a finite packet defeats all roots of orders $2,4,\ldots,2^{371}$, while some order at most
$2^{372}$ always succeeds. The comparisons use exact fractions; for $m<e_h$ written as
$m=(p-1)t+s$ with $0\le s<p-1$ and $t=\sum_jd_jp^j$ in base $p$, the numerator of $U_h(m)$
is $(p-1)\sum_j(j+1)p^jd_j$, which makes the threshold $m_h(\delta)$ accessible digit by
digit. The recorded outputs are:
\costfiles{
\item \artifact{grh-torsion-and-zero-data/verification/selection_figure.txt}
\item \artifact{grh-torsion-and-zero-data/verification/check_threshold.txt}
\item \artifact{grh-torsion-and-prime-sampling/verification.txt}}
\end{example}

\subsection{Packets of fixed cardinality}\label{cost:sec:cardinality}
{\raggedright\sources{grh_fixed_cardinality_torsion}}

Theorems~\ref{cost:thm:minimax} and~\ref{cost:thm:packet} are the two extremes of the packet
problem in setting (C). A packet with one member $(g,t)$ blocks only if $g$ blocks in the
aggregate sense, since $t\ge G_\rho(g)$, and the extremal polynomial of
Theorem~\ref{cost:thm:minimax} with threshold $G_\rho$ is such a packet; at the other end,
$H-h_0+1$ members suffice for Theorem~\ref{cost:thm:packet}. Between them the cost grows
like an iterated logarithm whose height is the number of members.

We stay in setting (C), with packets over $\OO$ as in Definition~\ref{cost:def:packet}, a
fixed $h_0\ge1$ and $\rho\ge r_{h_0}$, and we write $s_{h_0}=p^{1-h_0}$ as in
Theorem~\ref{cost:thm:transition}. For $n\ge1$, $H\ge h_0$ and $\delta\ge0$ let
$C_n(H,\delta)$ be the infimum of the costs of packets over $\OO$ with at most $n$ members
that block the shells $h_0,\dots,H$ at tolerance $\delta$. Let $H_n(B,\delta)$ be the least
integer $H\ge h_0$ such that every packet over $\OO$ with at most $n$ members and cost at
most $B$ has a point $z\in Z_{h_0}\cup\dots\cup Z_H$ with $\mathcal L(z)\le\delta$, and put
$H_n(B)=H_n(B,0)$. At tolerance zero such a point satisfies $v(g_i(z))\le t_i$ for every
member simultaneously. Unused members may be padded with the constant $1$ at threshold zero.
We write $\log_p^{\circ j}$ for the $j$-fold iterate of $\log_p$, the zeroth iterate being the
identity.

\begin{theorem}[Fixed-cardinality tower law; \stU]\label{cost:thm:cardinality}
For fixed $p$, $h_0$, $n$ and every $B\ge0$, $H_n(B)$ is finite, uniformly in
$\rho\ge r_{h_0}$. For $n=1$,
\[
 H_1(B)=h_0+\floor{B/s_{h_0}} .
\]
For fixed $n\ge2$, as $B\to\infty$,
\begin{equation}\label{cost:eq:tower}
 \frac B{s_{h_0}}-O_{p,h_0,n}(1)\ \le\ \log_p^{\circ(n-1)}H_n(B)\ \le\
 \frac B{s_{h_0}}+O_{p,h_0,n}\bigl(\log(2+B)\bigr).
\end{equation}
Equivalently, with $X_H=\log_p^{\circ(n-1)}(H+1)$, as $H\to\infty$,
\begin{equation}\label{cost:eq:towerinverse}
 s_{h_0}X_H-O_{p,h_0,n}\bigl(\log(2+X_H)\bigr)\le C_n(H,0)\le s_{h_0}X_H+O_{p,h_0,n}(1).
\end{equation}
The bounds are independent of the degrees of the polynomials, of the ramification of their
coefficients and of $\rho\ge r_{h_0}$.
\end{theorem}

For $p=2$ and $h_0=1$, one member requires budget $H$ to block the orders
$2,4,\dots,2^H$, two members require $\log_2H+O(\log\log H)$, and three require
$\log_2\log_2H+O(\log\log\log H)$, while with unrestricted cardinality the budget
$\sum_{h=h_0}^Hp^{1-h}$ of \eqref{cost:eq:zerotol} stays bounded. These statements concern
the order exponent $h$; the order of the torsion point is $p^h$.

\begin{theorem}[Exact finite optimization; \stU]\label{cost:thm:cardfinite}
For all $n\ge1$, $H\ge h_0$ and $\delta\ge0$ the infimum $C_n(H,\delta)$ is a minimum,
independent of $\rho\ge r_{h_0}$. It is attained by monic polynomials whose roots are torsion
points $\xi_h^b-1$ on the shells $h_0,\dots,H$, with at most $e_h$ roots of each member on
shell $h$, counted with multiplicity, and with thresholds $t_i=G_\rho(g_i)=v(g_i(0))$. It is
computed by the finite integer-state recurrence \eqref{cost:eq:cardstep} below, and
\[
 C_n(H,\delta)=C_n(H,0)\qquad(0\le\delta<r_H).
\]
\end{theorem}

For $n\ge H-h_0+1$ the minimum is the cost $\sum_{h=h_0}^Hc_h(\delta)$ of
Theorem~\ref{cost:thm:packet}, and for $n=1$ at tolerance zero it is the cost
$(H-h_0+1)s_{h_0}$ of Theorem~\ref{cost:thm:minimax}.

\paragraph{Root accounting for a packet.}
For a packet $(g_i,t_i)_i$ over $\OO$ and $h_0\le h\le H$, let $m_{i,h}$ be the number of
roots of $g_i$ of valuation $r_h$, counted with multiplicity, and put
\[
 w_{i,h}=m_{i,h}r_h,\qquad W_i(h)=\sum_{k=h_0}^hw_{i,k},\qquad
 S_i(h)=\sum_{h_0\le k<h}m_{i,k}(r_k-r_h)=\sum_{h_0\le k<h}w_{i,k}\bigl(1-p^{k-h}\bigr).
\]
For $z\in Z_h$ let $E_{i,h}(z)=\sum_\alpha(v(z-\alpha)-r_h)$, the sum over the roots of $g_i$
of valuation $r_h$. The \emph{activation level} of a member with $W_i(H)>0$ is the least $h$
with $m_{i,h}>0$.

\begin{lemma}[Packet root model]\label{cost:lem:cardmodel}
\leavevmode
\begin{enumerate}[label=\textup{(\roman*)}]
\item $W_i(H)\le\Delta_\rho(g_i)\le t_i$, and $S_i(h)\le W_i(h-1)$.
\item For $z\in Z_h$, $v(g_i(z))-t_i\le E_{i,h}(z)-S_i(h)$, with equality when
$t_i=G_\rho(g_i)$ and every root of $g_i$ has one of the valuations
$r_{h_0},\dots,r_H$.
\item A member without roots of valuation $r_h$ has no loss on $Z_h$. In particular a
blocked shell carries a root of some member.
\end{enumerate}
\end{lemma}

\begin{proof}
(i) is \eqref{cost:eq:budget} together with $G_0(g_i)\ge0$. For (ii), Lemma~\ref{cost:lem:roots}(ii)
gives $v(g_i(z))-G_{r_h}(g_i)=E_{i,h}(z)$, and $G_\rho(g_i)-G_{r_h}(g_i)\ge S_i(h)$ by the
argument for \eqref{cost:eq:slack}: by \eqref{cost:eq:gauss} every root contributes a
nonnegative amount, a root of valuation $r_k$, $k<h$, contributes $r_k-r_h$, and a root of
valuation $r_k$, $k\ge h$, contributes zero. Finally $t_i\ge G_\rho(g_i)$. Part (iii)
follows, since then $E_{i,h}=0$ and $S_i(h)\ge0$.
\end{proof}

\begin{proof}[Reduction in Theorem~\ref{cost:thm:cardfinite}]
Given a blocking packet, replace each root $\alpha$ of valuation $r_h$, $h_0\le h\le H$, by a
point $\beta$ of $Z_h$ maximizing $v(\alpha-\beta)$; as in the proof of
Lemma~\ref{cost:lem:shell}(i), $v(z-\alpha)\le v(z-\beta)$ for every $z\in Z_h$, so the
on-shell excess can only increase, and the contributions on other shells depend only on the
valuation. Remove all other roots and the leading coefficient, and reset each threshold to
the Gauss valuation of the new monic polynomial. By \eqref{cost:eq:gauss}, a removed root of
valuation different from all $r_h$ contributed $\min\set{r_h,v(\alpha)}-\min\set{\rho,v(\alpha)}\le0$
to $v(g_i(z))-G_\rho(g_i)$ at $z\in Z_h$, so no positive part decreases. The new threshold
is $\sum_hm_{i,h}r_h=W_i(H)\le t_i$ by Lemma~\ref{cost:lem:cardmodel}(i), and it equals
$v(g_i(0))$ because all roots have valuation at most $r_{h_0}\le\rho$. If a member has more
than $e_h$ roots on shell $h$, replace them by the $e_h$ points of $Z_h$, once each: its loss
on $Z_h$ becomes infinite, on other shells each removed root had a nonpositive contribution,
and the cost decreases. The reduced packets form a finite set, independent of $\rho$, and no
reduction increases the cost or destroys blocking. This proves the attainment and the
independence of $\rho$.
\end{proof}

\paragraph{The finite recurrence.}
Identify $Z_h$ with the unit residues $a$ modulo $p^h$ through $a\mapsto\xi_h^a-1$. In a
reduced packet, member $i$ places on shell $h$ a multiset of $m_i\le e_h$ marks $b$, and by
\eqref{cost:eq:distance} its on-shell excess at the leaf $a$, in units of $r_h$, is
\[
 V_i(a)=\sum_{b\ \mathrm{marked}}\bigl(p^{v_p(a-b)}-1\bigr),
\]
interpreted as $+\infty$ when $a$ is marked; for $m_i=e_h$ the whole shell may be used. The
state after shell $h-1$ is the list of pairs $(M_i,K_i)$, where $M_i$ is the number of
earlier roots of member $i$ and $K_i$ their total valuation in units of $r_{h-1}$; before
the first shell both are zero. Since $r_{h-1}=pr_h$, the slack $S_i(h)$ is $pK_i-M_i$ in
units of $r_h$, and no later root changes it. By Lemma~\ref{cost:lem:cardmodel}(ii), with
equality for reduced packets, a choice of layouts on shell $h$ is feasible exactly when
\begin{equation}\label{cost:eq:cardstep}
 \sum_i\max\set{0,\,V_i(a)-(pK_i-M_i)}>\delta e_h\qquad\text{for every unit }a\bmod p^h,
\end{equation}
and the state then passes to $(M_i+m_i,\ pK_i+m_i)$.
All finite terms are integers, so the terms may be clipped at $\floor{\delta e_h}+1$. The
recurrence keeps all reachable states, which may be identified up to permutation of the
members, and $C_n(H,\delta)$ is the minimum of $r_H\sum_iK_i$ over the final states. Every
feasible path is realized by the products of its torsion linear factors, so the recurrence
neither relaxes nor tightens the problem.

\begin{proof}[Stabilization in Theorem~\ref{cost:thm:cardfinite}]
In a reduced packet every finite loss on $Z_h$ is an integer multiple of $r_h$, by
\eqref{cost:eq:cardstep}. Hence a reduced packet blocking at tolerance zero has loss at
least $r_h\ge r_H$, or infinite loss, at every point of every shell $h\le H$, and it blocks
at every tolerance $\delta<r_H$. Monotonicity in $\delta$ gives the equality.
\end{proof}

\begin{corollary}[No inverse-tolerance singularity; \stU]\label{cost:cor:cardtol}
For fixed $p$, $h_0$, $n$ and $B$, $H_n(B,\delta)=H_n(B)$ for every
$0\le\delta<r_{H_n(B)}$. In particular the least guaranteed exponent stays bounded as
$\delta\downarrow0$, in contrast with Theorem~\ref{cost:thm:transition}\textup{(ii),(iii)} for
unrestricted packets.
\end{corollary}

\begin{proof}
Write $J=H_n(B)$, which is finite by Theorem~\ref{cost:thm:cardinality}. If $J=h_0$ the claim
is immediate. Otherwise a packet of cost at most $B$ blocks the shells through $J-1$ at
tolerance zero, and by attainment and stabilization it also blocks them at every tolerance
$\delta<r_J<r_{J-1}$. Thus $H_n(B,\delta)\ge J$, and the reverse inequality is monotonicity.
\end{proof}

\paragraph{A capped root average.}
For the upper bound put $B'=B+1$ and assume $B\ge2$; a smaller budget may be replaced by
$2$.

\begin{lemma}[Capped average over a shell]\label{cost:lem:cardcap}
If $v(\alpha)=r_h$, then
\[
 \frac1{e_h}\sum_{z\in Z_h}\min\set{v(z-\alpha)-r_h,\ B'}\le(h+B')\,r_h .
\]
\end{lemma}

\begin{proof}
Let $\beta\in Z_h$ be a nearest point to $\alpha$, so that $v(z-\alpha)\le v(z-\beta)$ on
$Z_h$. The monic polynomial $P_h(T)=\Phi_{p^h}(1+T)$ has the simple roots $Z_h$, and
differentiating $(X^{p^h}-1)/(X^{p^{h-1}}-1)$ at $X=1+\beta$ shows
$v(P_h'(\beta))=h-1/(p-1)$. Hence
$\sum_{z\ne\beta}(v(z-\beta)-r_h)=h-1/(p-1)-(e_h-1)r_h\le h$. The term at $\beta$ is at
most $B'$. Divide by $e_h$.
\end{proof}

\begin{lemma}[Forced growth on a blocked shell]\label{cost:lem:cardforced}
Let a packet of cost at most $B$ block shell $h$ at tolerance zero, and suppose that every
member with a root of valuation $r_h$ has $W_i(h-1)>0$. Then, summing over the members with
$W_i(h-1)>0$,
\begin{equation}\label{cost:eq:forced}
 \sum_i\frac{w_{i,h}}{S_i(h)}>\frac1{h+B'} .
\end{equation}
\end{lemma}

\begin{proof}
For these members $0<S_i(h)\le W_i(h-1)\le B$. At $z\in Z_h$ some member has
$v(g_i(z))>t_i$, so $E_{i,h}(z)>S_i(h)$ by Lemma~\ref{cost:lem:cardmodel}(ii), and this member
has a root of valuation $r_h$. Capping each root excess at $B'>B$ keeps this strict
inequality: either no term changes, or one capped term alone exceeds $S_i(h)$. Hence the sum
over the members with $W_i(h-1)>0$ of the capped excess divided by $S_i(h)$ exceeds one at
every point of $Z_h$. Average over $Z_h$ and apply Lemma~\ref{cost:lem:cardcap} to each root.
\end{proof}

\begin{lemma}[Delayed product potential]\label{cost:lem:cardinterval}
Let a packet of cost at most $B$ block every shell $a,a+1,\dots,b$ at tolerance zero, where
$h_0\le a\le b\le H$, and suppose the members with $W_i(b)>0$ are exactly the $k$ members with
$W_i(a)>0$, with activation levels $a_i\le a$. For an integer $d\ge1$ and a real $K>1$ put
\[
 \lambda=\frac{1-p^{-d}}K,\qquad L=\sum_{i=1}^k\log\bigl(Be_{a_i}\bigr),\qquad
 D=d+\ceil*{\frac{dL}{\log K}} .
\]
Then
\begin{equation}\label{cost:eq:interval}
 b+B'+1\le(a+B'+1+D)\exp(L/\lambda).
\end{equation}
\end{lemma}

\begin{proof}
Logarithms are natural. For $h\ge a$ let $P_h=\prod_{i=1}^kW_i(h)>0$ and, for $h>a$,
$R_h=P_h/P_{h-1}\ge1$. Since $W_i(a)\ge r_{a_i}$ and $W_i(b)\le B$,
$\sum_{h=a+1}^b\log R_h=\log(P_b/P_a)\le L$.

Discard the first $d$ shells after $a$. Call a later shell $h\le b$ irregular if
$P_{h-1}/P_{h-d}>K$. Each increment $\log R_j$ occurs in at most $d$ of the quotients
$\log(P_{h-1}/P_{h-d})$, so there are at most $\ceil{dL/\log K}$ irregular shells. On a
regular shell each ratio $W_i(h-1)/W_i(h-d)\ge1$ is at most $K$, because their product is,
and retaining only the roots of levels at most $h-d$ gives
\[
 S_i(h)\ge(1-p^{-d})\,W_i(h-d)\ge\lambda\,W_i(h-1).
\]
Every member with a root on shell $h$ is one of the $k$ members, so
Lemma~\ref{cost:lem:cardforced} applies, and
\[
 R_h=\prod_i\Bigl(1+\frac{w_{i,h}}{W_i(h-1)}\Bigr)\ge1+\lambda\sum_i\frac{w_{i,h}}{S_i(h)}
 >1+\frac\lambda{h+B'} .
\]
Since $0<\lambda<1$, concavity gives $\log(1+\lambda/x)\ge\lambda\log(1+1/x)$ for $x>0$.
The summands $\log(1+1/(h+B'))$ decrease in $h$, and at most $D$ of the shells
$a+1,\dots,b$ are not regular, so
\[
 L\ge\lambda\sum_{h\ \mathrm{regular}}\log\Bigl(1+\frac1{h+B'}\Bigr)
 \ge\lambda\log\frac{b+B'+1}{a+B'+1+D};
\]
if $D>b-a$ the right side is not positive. Exponentiate.
\end{proof}

The delay $d$ is essential to the leading constant: the elementary bound
$S_i(h)\ge(1-1/p)W_i(h-1)$ alone would put a factor $p/(p-1)$ into the exponent at each
activation.

\begin{proof}[Proof of the upper bound in Theorem~\ref{cost:thm:cardinality}]
Let a packet with at most $n$ members and cost at most $B$ block $h_0,\dots,H$ at tolerance
zero. Let $a_1\le\dots\le a_m$, $m\le n$, be the activation levels of the members with
$W_i(H)>0$, and put $a_{m+1}=H+1$. By Lemma~\ref{cost:lem:cardmodel}(iii), $a_1=h_0$, and the
shells before $a_2$ are blocked by the first member alone, which therefore blocks them in
the aggregate sense. Theorem~\ref{cost:thm:minimax}(i) and $\Delta_\rho(g_1)\le B$ give
\begin{equation}\label{cost:eq:a2}
 a_2\le h_0+B/s_{h_0},
\end{equation}
and if $m=1$ the same bound holds for $H+1$. For $k\ge2$, on a nonempty interval
$[a_k,a_{k+1}-1]$ exactly the first $k$ members are active, and Lemma~\ref{cost:lem:cardinterval}
applies. Put
\[
 d=\ceil{\log_pB},\quad K=1+B^{-1},\quad\lambda=\frac{1-p^{-d}}K,\qquad U_1=h_0,\quad
 U_2=h_0+\frac B{s_{h_0}},
\]
and for $k=2,\dots,n$
\[
 L_k=k\log\bigl(B(p-1)\bigr)+\log p\sum_{i=1}^k(U_i-1),\qquad
 U_{k+1}=\Bigl(U_k+B'+1+d+\ceil*{\frac{dL_k}{\log K}}\Bigr)\exp(L_k/\lambda)+1 .
\]
The sequence $U_k$ increases, and induction using \eqref{cost:eq:interval} and
\eqref{cost:eq:a2} gives $a_i\le U_i$ for $i\le m+1$, since $\log(Be_{a_i})=\log(B(p-1))+(a_i-1)\log p$.
Hence $H+1\le U_{n+1}$ for every blocker, and $H_n(B)\le\ceil{U_{n+1}}$. This proves
finiteness at tolerance zero, uniformly in $\rho$, without compactness or a limit in the
tolerance.

As $B\to\infty$, $d=O_p(\log B)$, $\lambda^{-1}=1+O(B^{-1})$ and $(\log K)^{-1}=O(B)$. For
$k=2$, $L_2/\log p=(U_2-1)+(h_0-1)+2\log_p(B(p-1))=B/s_{h_0}+O_{p,h_0}(\log B)$, and the
logarithm of the prefactor of $\exp(L_2/\lambda)$ is $O_{p,h_0}(\log B)$, so
$\log_pU_3=B/s_{h_0}+O_{p,h_0}(\log B)$. For fixed $k\ge3$, $L_k=O_{p,h_0,n}(U_k)$ and
$U_k\ge B$, and the logarithm of the prefactor is $O_{p,h_0,n}(\log U_k)$; thus
$\log_pU_{k+1}\le c_1U_k$ for a constant $c_1\ge1$. Since
$\log_p^{\circ j}(c_1x)\le\log_p^{\circ j}(x)+O(1)$ for $j\ge1$ and large $x$, iteration
gives $\log_p^{\circ(n-1)}U_{n+1}\le B/s_{h_0}+O_{p,h_0,n}(\log B)$, and the same bound holds
for $H_n(B)$. For $n=1$, \eqref{cost:eq:a2} with $H+1$ in place of $a_2$ gives
$H_1(B)\le h_0+\floor{B/s_{h_0}}$, and equality is Corollary~\ref{cost:cor:endpoint}.
\end{proof}

\begin{proof}[Proof of the lower bound and of \eqref{cost:eq:towerinverse}]
For large $B$ put $a_1=h_0$, $\ell_1=\floor{(B-(n-1))/s_{h_0}}$, $a_2=h_0+\ell_1$, and for
$2\le i\le n$ let $\ell_i=p^{a_i-1}$ and $a_{i+1}=a_i+\ell_i$. On the block of shells
$a_i,\dots,a_{i+1}-1$ use the extremal polynomial of Theorem~\ref{cost:thm:minimax}(i)
with $a_i$ in place of $h_0$, with threshold its Gauss valuation; its cost is $s_{h_0}\ell_1$
for $i=1$ and $\ell_ip^{1-a_i}=1$ for $i\ge2$. Each member blocks its own block whatever the
others do, and the total cost is at most $B$. Hence $H_n(B)\ge a_{n+1}$. Since
$\log_p(a_i+p^{a_i-1})=a_i-1+o(1)$ as $a_i\to\infty$, $n-1$ iterated logarithms of
$a_{n+1}$ give $a_2+O_{p,n}(1)=B/s_{h_0}+O_{p,h_0,n}(1)$. The construction uses
coefficients in one cyclotomic field, depending on $H$, and the common witness $0$.

By attainment, $C_n(H,0)\le B$ holds exactly when $H_n(B)>H$. Choosing
$B=s_{h_0}(X_H+c)$ with $c$ large, the lower bound in \eqref{cost:eq:tower} gives
$H_n(B)>H$, so $C_n(H,0)\le s_{h_0}X_H+O(1)$. Substituting $B=C_n(H,0)=O(X_H)$ into the
upper bound of \eqref{cost:eq:tower} gives the other inequality.
\end{proof}

\begin{example}[Small exact values]\label{cost:ex:cardinality}
The recurrence \eqref{cost:eq:cardstep}, run exactly, gives the following costs.
\begin{center}
\begin{tabular}{rrrrrrl}
\toprule
$p$ & $h_0$ & $H$ & $n$ & $\delta$ & $C_n(H,\delta)$ & \\
\midrule
2&1&3&1&0&3&one member\\
2&1&3&2&0&2&\\
2&1&3&3&0&$7/4$&unrestricted optimum\\
3&1&2&2&0&$4/3$&\\
2&2&3&1&0&1&\\
2&2&3&2&0&$3/4$&\\
2&1&3&1&$1/4$&3&\\
2&1&3&2&$1/4$&2&\\
2&1&3&3&$1/4$&2&\\
\bottomrule
\end{tabular}
\end{center}
The one-member values agree with Theorem~\ref{cost:thm:minimax}, and for $n\ge H-h_0+1$ the
values are those of Theorem~\ref{cost:thm:packet}. The small zero-tolerance optima agree
with the best partition into consecutive blocks with one member each; the recurrence
also allows mixed assignments within a shell, and it is not known whether consecutive
blocks are always optimal. The same script checks the capacity formula of
Lemma~\ref{cost:lem:shell} for $27$ pairs of shell and root count.
\costfiles{
\item \artifact{grh_fixed_cardinality_torsion/packet_exact.py}
\item \artifact{grh_fixed_cardinality_torsion/verification.txt}}
\end{example}

\paragraph{Exact selection in any dimension.}
For a nonzero integral power series $F=\sum_\nu c_\nu T^\nu\in\OO[[T_1,\dots,T_d]]$ put
$G_\rho(F)=\min_\nu\bigl(v(c_\nu)+\rho\abs\nu\bigr)$.

\begin{theorem}[Common tuple at exact Gauss values; \stU]\label{cost:thm:cardmulti}
Let $F_1,\dots,F_n\in\OO[[T_1,\dots,T_d]]$ be nonzero, $d\ge1$, put $A_i=G_\rho(F_i)$ and
suppose $\sum_iA_i\le B$. Then there are $z_1,\dots,z_d\in Z_{h_0}\cup\dots\cup Z_{H_n(B)}$
with
\[
 v\bigl(F_i(z_1,\dots,z_d)\bigr)\le A_i\qquad(1\le i\le n).
\]
The least endpoint with this property is $H_n(B)$, for every $d$. If instead witnesses
$x_i\in\Cp^d$ with coordinates of valuation at least $\rho$ satisfy
$\sum_iv(F_i(x_i))\le B$, the same endpoint works with $v(F_i(x_i))$ in place of $A_i$.
\end{theorem}

\begin{proof}
First let the $F_i$ be polynomials. Regard $F_i$ as a polynomial in $T_2,\dots,T_d$ with
coefficients $g_{i,\beta}(T_1)$, and choose $\beta_i$ with
$G_\rho(g_{i,\beta_i})+\rho\abs{\beta_i}=A_i$. The polynomials $g_{i,\beta_i}$ are integral
and nonzero, and their Gauss valuations add up to at most $B$. By the definition of
$H_n(B)$, some $z_1$ in the shells $h_0,\dots,H_n(B)$ satisfies
$v(g_{i,\beta_i}(z_1))\le G_\rho(g_{i,\beta_i})$ for every $i$. Then $F_i(z_1,\cdot)$ is
nonzero and integral, and its Gauss valuation is at most $A_i$, since its coefficient at
$\beta_i$ already gives this bound. Iterate through the remaining coordinates; the choice of
$\beta_i$ may change from step to step, and no loss accumulates with the dimension.

For power series, $\rho>0$ and the integrality of the coefficients imply that $G_\rho(F_i)$
is attained in finite total degree. Truncate at a total degree $D$ with
$Dr_{H_n(B)}>\max_iA_i$. The truncations have the same $A_i$, and the omitted terms have
valuation greater than $A_i$ at every tuple of the tested shells, so the ultrametric
inequality transfers the conclusion. Convergent witness evaluations satisfy
$v(F_i(x_i))\ge A_i$, which gives the last assertion. Sharpness follows from an extremal
one-variable packet for the shells $h_0,\dots,H_n(B)-1$, with thresholds equal to the Gauss
valuations and witness $0$ (Theorem~\ref{cost:thm:cardfinite}), independent of the other
coordinates.
\end{proof}

With $n$ fixed this is a zero-tolerance counterpart of \cref{tors:thm:packet}: the endpoint
does not depend on the dimension, the degrees or the coefficient field.

\subsection{Separate bounds with a positive tolerance}\label{tors:sec:tolerance}
\sources{grh_quantitative_twists_2026-11-25; grh_exact_torsion_transfer_2026-11-25; grh_method_transfers_2026-11-25_1; grh_beyond_dirichlet_grh_2026-11-25}

We now ask for a common torsion tuple at which \emph{each} member of a family is within a
prescribed tolerance of its own target. Two regimes have to be distinguished. For a family
of fixed cardinality $n$ the exponent can be taken polynomial in the inverse tolerances,
and a positive fraction of the primitive roots of one order is good
(\cref{tors:thm:fraction,tors:thm:packet}); for $n=1$ the dependence is linear
(\cref{tors:prop:linear}). For families of arbitrary cardinality with a finite total
budget the harmonic argument of the source gives an exponent that is exponential in
budget over tolerance (\cref{tors:thm:uniform}), and this rate cannot be improved
(\cref{tors:rem:tolerance}). Throughout, $c=1-1/p$ and $q=1/p$.

\begin{lemma}[Capped shell average]\label{tors:lem:capped}
Let $h\ge1$, let $v(\alpha)=r_h$ and $C>0$, and put
$E_\alpha(z)=v(z-\alpha)-r_h\in[0,+\infty]$ for $z\in Z_h$. Then
\[
  \frac1{e_h}\sum_{z\in Z_h}\min\{E_\alpha(z),C\}\le r_h\,(h+C).
\]
\end{lemma}

\begin{proof}
Choose $z_*=\zeta_*-1\in Z_h$ maximizing $v(z-\alpha)$. For $z=\zeta-1\ne z_*$ one has
$v(z_*-\alpha)\ge v(z-\alpha)$, so the ultrametric inequality gives
$v(z-\alpha)\le v(z-z_*)=v(\zeta-\zeta_*)$. The summand at $z_*$ is at most $C$, and every other summand is
at most $v(\zeta-\zeta_*)$. Since
$\prod_{\zeta\ne\zeta_*}(\zeta_*-\zeta)=\Phi'_{p^h}(\zeta_*)$, the sum is at most
$C+v(\Phi'_{p^h}(\zeta_*))$. Differentiating
$\Phi_{p^h}(X)\,(X^{p^{h-1}}-1)=X^{p^h}-1$ at $\zeta_*$ gives
$v(\Phi'_{p^h}(\zeta_*))=h-1/(p-1)\le h$, because $\zeta_*^{p^{h-1}}$ is a primitive
$p$-th root of unity. Dividing by $e_h$ proves the lemma; the case $p=2$, $h=1$, in which
the shell has one point, is included.
\end{proof}

\begin{theorem}[Polynomial exponent and a positive fraction of good roots; \stU]
\label{tors:thm:fraction}
Let $g_1,\dots,g_n\in\overline{\Z}_p[T]$ be nonzero, $n\ge1$, let $\rho>0$, and let
$B_i\ge0$ and
$\delta_i>0$ be such that
\[
  A_i:=G_\rho(g_i)\le B_i\qquad(1\le i\le n).
\]
Let $h_0\ge1$ satisfy $r_{h_0}\le\rho$, let $0<\eta\le1$, and define
\begin{equation}\label{tors:eq:CP}
  C=1+\max_i(B_i+\delta_i),\qquad P=\prod_{i=1}^n\Bigl(1+\frac{cB_i}{\delta_i}\Bigr).
\end{equation}
Let $H\ge h_0$ be an integer with
\begin{equation}\label{tors:eq:Hcondition}
  \Bigl(\frac{H+C+1}{h_0+C}\Bigr)^{c\eta}\ge P ;
\end{equation}
one admissible choice is $H=\ceil[\big]{(h_0+C)P^{1/(c\eta)}}$. Then there is $h$ with
$h_0\le h\le H$ such that strictly more than a fraction $1-\eta$ of the points
$z\in Z_h$ satisfy all the inequalities
\begin{equation}\label{tors:eq:desired}
  v(g_i(z))\le A_i+\delta_i\qquad(1\le i\le n).
\end{equation}
In particular all these values are nonzero. The bound does not depend on degrees, roots or
coefficient fields.
\end{theorem}

For $\eta=1$ the assertion is existence; for $\eta=\tfrac12$ more than half of the
primitive roots of one order are good. The comparison is with the fixed $G_\rho$, not with
the moving $G_{r_h}$. Exact rational checks of the inequalities of the following proof
on small shells are recorded in:
\torsfiles{
\item \artifact{grh_quantitative_twists_2026-11-25/verification.txt}}

\begin{proof}
For each $i$ and $h\ge h_0$ let $R_{ih}$ be the multiset of roots of $g_i$ of valuation
$r_h$, and put
\[
  w_{ih}=r_h\,\#R_{ih},\qquad W_{i,h}=\sum_{k=h_0}^hw_{ik},\qquad W_{i,h_0-1}=0 .
\]
By \cref{tors:lem:shells}(a),
\begin{equation}\label{tors:eq:total}
  W_{i,H}\le A_i\le B_i .
\end{equation}
Let $s_{ih}=A_i-G_{r_h}(g_i)$ be the slack at shell $h$. A root on an earlier shell
$k<h$ contributes $r_k-r_h\ge c\,r_k$ to it, so \eqref{tors:eq:slack} gives
\begin{equation}\label{tors:eq:slackc}
  0\le c\,W_{i,h-1}\le s_{ih}\le B_i .
\end{equation}
Let $D_{ih}(z)=\sum_{\alpha\in R_{ih}}\min\{E_\alpha(z),C\}$. If \eqref{tors:eq:desired}
fails for $i$ at $z\in Z_h$, then by \eqref{tors:eq:value} the sum of the uncapped
excesses $E_\alpha(z)$, $\alpha\in R_{ih}$, exceeds $s_{ih}+\delta_i$. Since
$C>B_i+\delta_i\ge s_{ih}+\delta_i$, the capped sum still exceeds it: either one term
reaches the cap or no term changes. Hence
\[
  \frac{D_{ih}(z)}{\delta_i+cW_{i,h-1}}>1\qquad\text{whenever constraint $i$ fails at }z.
\]
Call $z$ bad if it violates at least one constraint. If at least a fraction $\eta$ of
$Z_h$ is bad, the average over $Z_h$ of the sum over $i$ of these nonnegative ratios is
strictly greater than $\eta$; the inequality is strict also for $\eta=1$, because the
ratio sum is greater than one at every bad point. \Cref{tors:lem:capped} bounds the
average of $D_{ih}$ by $w_{ih}(h+C)$, so
\begin{equation}\label{tors:eq:bad}
  \sum_{i=1}^n\frac{w_{ih}}{\delta_i+cW_{i,h-1}}>\frac{\eta}{h+C}.
\end{equation}
Set $\mathcal P_h=\prod_i(1+cW_{i,h}/\delta_i)$. If the bad fraction of shell $h$ is at
least $\eta$, then
\[
  \frac{\mathcal P_h}{\mathcal P_{h-1}}
  =\prod_i\Bigl(1+\frac{c\,w_{ih}}{\delta_i+cW_{i,h-1}}\Bigr)
  \ge1+\sum_i\frac{c\,w_{ih}}{\delta_i+cW_{i,h-1}}>1+\frac{c\eta}{h+C}.
\]
Suppose that all shells $h_0,\dots,H$ have bad fraction at least $\eta$. Since
$\mathcal P_{h_0-1}=1$ and $0<c\eta\le1$, concavity of $u\mapsto u^{c\eta}$ gives
\[
  \mathcal P_H>\prod_{h=h_0}^H\Bigl(1+\frac{c\eta}{h+C}\Bigr)
  \ge\Bigl(\prod_{h=h_0}^H\frac{h+C+1}{h+C}\Bigr)^{c\eta}
  =\Bigl(\frac{H+C+1}{h_0+C}\Bigr)^{c\eta}.
\]
But \eqref{tors:eq:total} implies $\mathcal P_H\le P$, contradicting
\eqref{tors:eq:Hcondition}. Hence some shell has bad fraction less than $\eta$.
\end{proof}

For one polynomial the argument can be run with the exact saving
$r_k-r_h=(1-q^{h-k})r_k$ in place of its lower bound $c\,r_k$. This lowers the dependence
on the inverse tolerance from the exponent $1/c=p/(p-1)$ to the exponent one.

\begin{proposition}[Linear dependence on the inverse tolerance; \stU]\label{tors:prop:linear}
Let $\rho>0$, $B\ge0$, $\delta>0$, and let $h_0\ge1$ satisfy $r_{h_0}\le\rho$. Put
\[
  t_0=h_0+B+\delta+1,\qquad \kappa=1-\frac1{(p-1)(t_0+1)}>\frac12,\qquad
  H_+=h_0+\ceil*{\frac{B\,t_0}{\kappa\,\delta}} .
\]
For every nonzero $g\in\OO[T]$ with $G_\rho(g)\le B$ there are $h$ with
$h_0\le h\le H_+$ and $z\in Z_h$ such that $v(g(z))\le G_\rho(g)+\delta$.
\end{proposition}

\begin{proof}
Let $A=G_\rho(g)$, let $R_h$ be the multiset of roots of valuation $r_h$, and put
\[
  w_h=r_h\,\#R_h,\qquad W_h=\sum_{k=h_0}^hw_k,\qquad
  y_h=\sum_{k=h_0}^hq^{h-k}w_k,\qquad W_{h_0-1}=y_{h_0-1}=0 .
\]
As in \eqref{tors:eq:total}, $W_H\le A\le B$ for every $H$. By \eqref{tors:eq:slack}
the slack $s_h=A-G_{r_h}(g)$ satisfies
\begin{equation}\label{tors:eq:exactslack}
  s_h\ge\sum_{k=h_0}^{h-1}(r_k-r_h)\,\#R_k=W_{h-1}-q\,y_{h-1},
\end{equation}
because $r_h/r_k=q^{h-k}$. Suppose that every $z\in Z_h$ violates
$v(g(z))\le A+\delta$, and let $C=B+\delta+1$. By \eqref{tors:eq:value} the sum of the
excesses $E_\alpha(z)$, $\alpha\in R_h$, is greater than $s_h+\delta$ at every $z$, and
capping at $C$ preserves this, since $s_h+\delta\le B+\delta<C$. Averaging over $Z_h$ and
applying \cref{tors:lem:capped} and \eqref{tors:eq:exactslack} gives
\begin{equation}\label{tors:eq:failure}
  w_h>\frac{Q_h}{h+C},\qquad Q_h:=\delta+W_{h-1}-q\,y_{h-1}.
\end{equation}
The quantities satisfy
\begin{equation}\label{tors:eq:update}
  y_h=q\,y_{h-1}+w_h,\qquad Q_{h+1}=Q_h+c\,y_h .
\end{equation}
Suppose now that all shells $h_0,\dots,H$ fail. Define a comparison sequence by
$Q'_{h_0}=\delta$, $y'_{h_0-1}=0$ and
\[
  u_h=\frac{Q'_h}{t_h},\qquad y'_h=q\,y'_{h-1}+u_h,\qquad Q'_{h+1}=Q'_h+c\,y'_h,
  \qquad t_h=h+C .
\]
All coefficients in \eqref{tors:eq:update} are nonnegative, so \eqref{tors:eq:failure}
and induction give $Q_h\ge Q'_h$, $y_h\ge y'_h$ and $w_h>u_h$ for $h_0\le h\le H$. Put
$d_h=u_h-c\,y'_{h-1}$. Direct substitution yields
\begin{equation}\label{tors:eq:defect}
  u_{h+1}-u_h=-\frac{q\,d_h}{t_h+1},\qquad d_{h+1}=q\,\frac{t_h}{t_h+1}\,d_h .
\end{equation}
Here $t_{h_0}=t_0$ and $u_{h_0}=d_{h_0}=\delta/t_0$, so
\[
  d_{h_0+j}=\frac{\delta\,q^j}{t_0+j},\qquad
  u_{h_0+j}=\frac{\delta}{t_0}-\delta\sum_{k=0}^{j-1}\frac{q^{k+1}}{(t_0+k)(t_0+k+1)}
  \ \ge\ \frac{\kappa\,\delta}{t_0};
\]
the last inequality bounds each denominator below by $t_0(t_0+1)$ and sums the geometric
series. If every shell through $H_+$ failed, we would obtain
$B\ge W_{H_+}>(H_+-h_0+1)\kappa\delta/t_0>B$.
\end{proof}

\begin{remark}\label{tors:rem:linear}
For fixed $p,h_0,B$ and $0<\delta\le1$ one has $H_+=O(1/\delta)$, whereas
\cref{tors:thm:fraction} with $n=1$, $\eta=1$ gives
$H=O(\delta^{-p/(p-1)})$. The comparison in \cref{tors:prop:linear} is the one of
\cref{tors:thm:selection} with a tolerance added, so the endpoint
$\torsHs(B)=h_0+\floor{Bp^{h_0-1}}$ of that theorem, which needs no tolerance, is also
admissible; it is the smaller one unless $\delta$ is large. For example, for $p=3$,
$h_0=2$, $B=1$ and $\delta=\tfrac12,\tfrac1{10},\tfrac1{100},\tfrac1{1000}$ the explicit
choice in \cref{tors:thm:fraction} gives $17$, $88$, $2233$, $69026$ and
\cref{tors:prop:linear} gives $12$, $48$, $448$, $4448$
(\artifact{grh_exact_torsion_transfer_2026-11-25/data/bounds.csv}), while
$\torsHs(1)=5$. For $p=3$,
$h_0=3$, $B=1$, $\delta=1$ one finds $H_+=10<\torsHs(1)=12$. The least admissible
endpoint for one polynomial and a given tolerance is determined exactly, by a recursion
over the shells, in \cref{sec:cost}. Neither bound is polynomial for the order $p^H$
itself.
\end{remark}

\paragraph{Several variables and power series.}
The polynomial bound extends to several variables by replacing one coordinate at a time.
The positive fraction of good roots at each step gives many good tuples.

\begin{theorem}[Families of fixed cardinality; \stU]\label{tors:thm:packet}
Let $F_1,\dots,F_n\in\OO[[T_1,\dots,T_s]]$, and for each $i$ let $x_i\in\Cp^s$ be a
witness of depth $\rho>0$ with
\[
  t_i=v(F_i(x_i))\le B_i<\infty .
\]
Fix $s_0\ge s\ge1$, tolerances $\eps_i>0$, a number $0<\eta\le1$ and $h_0\ge1$ with
$r_{h_0}\le\rho$. Let $H$ satisfy \eqref{tors:eq:Hcondition}, where $C$ and $P$ are
defined by \eqref{tors:eq:CP} with the budgets $B_i+\eps_i$ and the tolerances
$\delta_i=\eps_i/s_0$. Then there are at least
\[
  K^s\ \text{distinct tuples }z\in\torsZ{h_0}{H}^s,\qquad
  K=\floor{(1-\eta)e_{h_0}}+1,
\]
such that
\[
  v(F_i(z))\le t_i+\eps_i\qquad\text{for all }i .
\]
\end{theorem}

\begin{proof}
First let all $F_i$ be polynomials with coefficients in $\overline{\Z}_p$ and all
witnesses lie in $\overline{\Q}_p^{\,s}$. Replace the coordinates one at a time. At a
common prefix $(z_1,\dots,z_{j-1})$ of replacements, freeze in each $F_i$ the coordinates
after the $j$-th at the witness $x_i$; this gives polynomials $g_i(T)$ in one variable
with coefficients in $\overline{\Z}_p$ and with $v(g_i(x_{ij}))\le t_i+(j-1)\delta_i$. Hence $g_i\ne0$ and, by
\eqref{tors:eq:witness}, $G_\rho(g_i)\le B_i+\eps_i$. \Cref{tors:thm:fraction} provides
a shell $h\in[h_0,H]$ on which more than $(1-\eta)e_h\ge(1-\eta)e_{h_0}$ points, hence at
least $K$ points, satisfy $v(g_i(z_j))\le G_\rho(g_i)+\delta_i\le t_i+j\delta_i$ for all
$i$. Keep all these branches. After $s$ steps every branch has loss at most
$s\delta_i\le\eps_i$, and distinct prefixes give distinct tuples. The bound $H$ was
chosen before any polynomial was seen.

For series over $\OO$ and witnesses in $\Cp^s$, choose $M>\max\{\rho,\max_i(B_i+\eps_i)\}$
and apply \cref{tors:lem:truncation} with $a=0$ and $r=r_H$; the coordinates of the
witnesses and of all points of the grid have valuation at least $r_H$. The truncations
$P_i$ satisfy $v(F_i(y)-P_i(y))>M$ at all these points. By \cref{tors:lem:approx} there
are polynomials $Q_i$ over $\overline{\Z}_p$ and witnesses $x'_i\in\overline{\Q}_p^{\,s}$
of depth $\rho$ with $v(Q_i(y)-P_i(y))>M$ on $\OO^s$ and $v(Q_i(x'_i)-P_i(x_i))>M$; hence
$v(Q_i(x'_i))=t_i$, since $t_i\le B_i<M$. The polynomial case gives $K^s$ tuples with
$v(Q_i(z))\le t_i+\eps_i<M$, and at each of them $v(F_i(z))=v(Q_i(z))$.
\end{proof}

To obtain a total positive excess at most $\eps$ for a family of $n$ members one can take
$\eps_i=\eps/n$; the exponent in \eqref{tors:eq:Hcondition} then depends on $n$. At zero
tolerance, \cref{cost:thm:cardmulti} gives the exact least endpoint for families of fixed
cardinality $n$. The next result is independent of the cardinality, at the price of an
exponential bound.

\begin{lemma}[One coordinate, harmonic choice of the order]\label{tors:lem:harmonic}
Let $\rho>0$, $0<\delta<1$ and $C_0\ge0$, and let $H$ satisfy
\[
  C_0<\delta\sum_{h=h_*(\rho)}^{H}\frac1{h+1}.
\]
If $g\in\overline{\Z}_p[T]$ is nonzero and $G_\rho(g)\le C_0$, there are $h$ with
$h_*(\rho)\le h\le H$ and $z\in Z_h$ such that
\[
  v(g(z))\le G_{r_h}(g)+\delta .
\]
\end{lemma}

\begin{proof}
Let $R_h$ be the multiset of roots of $g$ of valuation $r_h$. By
\cref{tors:lem:shells}(a),
\[
  \sum_{h=h_*(\rho)}^{H}r_h\,\#R_h\le G_\rho(g)\le C_0 ,
\]
so some $h$ in this range has $r_h\,\#R_h<\delta/(h+1)$. For $\alpha\in R_h$ the number
$1+\alpha$ is a unit, and
\[
  \sum_{z\in Z_h}\min\{v(z-\alpha),1\}\le h+1 .
\]
To see this, consider $Q(Y)=(Y+1+\alpha)^{p^h}-1$. Its coefficient of $Y$ has valuation
$h$, so $G_1(Q)\le h+1$; the numbers $z-\alpha$, $z\in Z_h$, are among its roots, and
\eqref{tors:eq:gauss} for the monic integral polynomial $Q$ gives the inequality. Sum
over $\alpha\in R_h$ and average over the $e_h=1/r_h$ points of $Z_h$: some $z\in Z_h$
satisfies
\[
  \sum_{\alpha\in R_h}\min\{v(z-\alpha),1\}\le r_h\,\#R_h\,(h+1)<\delta<1 .
\]
No summand equals $1$, so the minima can be removed and no difference vanishes. By
\eqref{tors:eq:value},
$v(g(z))=G_{r_h}(g)+\sum_{\alpha\in R_h}(v(z-\alpha)-r_h)\le G_{r_h}(g)+\delta$.
\end{proof}

Fix an upper bound $s_0\ge1$ for the number of variables, a depth $\rho>0$, a total
budget $B\ge0$ and an allowed loss $\eps>0$. Put
\begin{equation}\label{tors:eq:harmonicparam}
  h_*=h_*(\rho),\qquad \delta=\frac{\min\{\eps,1\}}{s_0+2},
\end{equation}
and let $H$ be an integer with
\begin{equation}\label{tors:eq:harmonic}
  B+1<\delta\sum_{h=h_*}^{H}\frac1{h+1}.
\end{equation}
An explicit admissible choice is
\begin{equation}\label{tors:eq:harmonicH}
  H=\ceil*{(h_*+1)\exp\Bigl(\frac{B+1}{\delta}\Bigr)},
\end{equation}
because $\sum_{h=h_*}^H(h+1)^{-1}\ge\log\bigl((H+2)/(h_*+1)\bigr)$. For nonnegative
numbers indexed by an arbitrary set, a sum means the supremum of the finite partial sums;
no countability is assumed.

\begin{theorem}[Uniform selection for families of bounded series; \stU]
\label{tors:thm:uniform}
Keep the parameters \eqref{tors:eq:harmonicparam}--\eqref{tors:eq:harmonic}. Let $I$ be
any index set and $1\le s\le s_0$. For $i\in I$ let $F_i\in\OO[[T_1,\dots,T_s]]$ and let
$x_i\in\Cp^s$ be a witness of depth $\rho$ with $t_i=v(F_i(x_i))$, and assume
\[
  \sum_{i\in I}t_i\le B .
\]
Then there is a single torsion tuple $z\in\torsZ{h_*}{H}^s$ such that
\begin{equation}\label{tors:eq:excess}
  \sum_{i\in I}\bigl(v(F_i(z))-t_i\bigr)_+\le\eps .
\end{equation}
Consequently every $F_i(z)$ is nonzero, $v(F_i(z))\le t_i+\eps$ for every $i$, and
$\sum_{i\in I}v(F_i(z))\le B+\eps$. The bound $H$ is independent of the cardinality of
$I$, of the complexity of the series, and of the fields generated by the coefficients and
the witnesses.
\end{theorem}

\begin{proof}
\emph{Finite families of polynomials.} Let $I$ be finite, all $F_i$ polynomials with
coefficients in $\overline{\Z}_p$ and all witnesses in $\overline{\Q}_p^{\,s}$. Replace
the coordinates one at a time, using the same replacement for all $F_i$. At a given step
free the coordinate to be replaced in each polynomial and freeze its other coordinates at
its current point; this gives polynomials $g_i$ in one variable with coefficients in
$\overline{\Z}_p$.
Let $c_i$ be the valuations of the current values. Their sum is at most $B+1$: after
$k\le s_0$ steps the total increase is at most $k\delta\le1$, as shown below. The
coordinate to be replaced has valuation at least $\rho$, so by \cref{tors:lem:gauss} and
\eqref{tors:eq:witness}
\[
  G_\rho\Bigl(\prod_ig_i\Bigr)=\sum_iG_\rho(g_i)\le\sum_ic_i\le B+1 .
\]
Apply \cref{tors:lem:harmonic} to $\prod_ig_i$ with $C_0=B+1$, and let
$a_i=v(g_i(z))$, $b_i=G_{r_h}(g_i)$ at the point $z\in Z_h$ it provides. Then
\[
  a_i\ge b_i,\qquad c_i\ge b_i,\qquad \sum_ia_i\le\sum_ib_i+\delta ;
\]
the inequality $c_i\ge G_\rho(g_i)\ge b_i$ uses $r_h\le\rho$. Therefore
\[
  \sum_i(a_i-c_i)_+\le\sum_i(a_i-b_i)\le\delta .
\]
In particular the total valuation grows by at most $\delta$ at each step, which justifies
the bound used above. The positive part of a sum of changes is at most the sum of their
positive parts, so after all $s$ coordinates the total positive excess is at most
$s\delta\le\eps$. Every selected coordinate lies in $\torsZ{h_*}{H}$.

\emph{Finite families of series.} Let $I$ be finite. Choose $M>B+\eps+\rho+1$ and apply
\cref{tors:lem:truncation} with $a=0$ and $r=r_H$; the coordinates of the witnesses and
of all points of the grid have valuation at least $r_H$, so the truncations $P_i$ satisfy
$v(F_i(y)-P_i(y))>M$ at all these points. By \cref{tors:lem:approx} there are
polynomials $Q_i$ over $\overline{\Z}_p$ and witnesses $x'_i\in\overline{\Q}_p^{\,s}$ of
depth $\rho$ with $v(Q_i(y)-P_i(y))>M$ on $\OO^s$ and $v(Q_i(x'_i)-F_i(x_i))>M$, hence
$v(Q_i(x'_i))=t_i$, since $t_i\le B<M$. The first step gives $z\in\torsZ{h_*}{H}^s$ with
total positive excess at most $\eps$ for the $Q_i$ and the $x'_i$. Each $v(Q_i(z))$ is at
most $t_i+\eps\le B+\eps<M$, so $v(F_i(z))=v(Q_i(z))$, and \eqref{tors:eq:excess} holds
for the $F_i$ without any further loss.

\emph{Arbitrary families.} For a finite $J\subseteq I$ let $S_J$ be the set of
$z\in\torsZ{h_*}{H}^s$ whose total positive excess over $J$ is at most $\eps$. Every $S_J$
is nonempty, and $S_{J_1\cup\dots\cup J_k}\subseteq S_{J_1}\cap\dots\cap S_{J_k}$ because
the summands are nonnegative. These subsets of a finite set have the finite intersection
property, so they have a common point $z$. Each singleton $J=\set{i}$ excludes an
infinite valuation, and summing $v(F_i(z))\le t_i+(v(F_i(z))-t_i)_+$ gives the last
bound.
\end{proof}

\begin{remark}[Comparison of the tolerance bounds]\label{tors:rem:tolerance}
(1) \Cref{tors:lem:harmonic} compares with the moving target $G_{r_h}(g)$ and is the
root-counting and averaging argument of the source; \cref{tors:thm:uniform} retains the
aggregate information that the final statement of the source discards, and extends it,
by truncation and algebraic approximation, to bounded series over $\OO$ and, by
compactness of the finite grid, to infinite families. The version for several powers,
with separate bounds for a finite family, is \cref{tors:prop:powers}(c).

(2) Neither of \cref{tors:thm:packet,tors:thm:uniform} contains the other. The exponent
in \cref{tors:thm:packet} is polynomial in the inverse tolerances for a fixed number $n$
of constraints, but its degree grows with $n$; the exponent \eqref{tors:eq:harmonicH} is
exponential in $(B+1)/\delta$ and independent of $n$. For equal budgets $B_i\le B_0$ and
tolerances $\delta_i=\delta$ the former is
$O\bigl((h_0+B_0+\delta+1)(1+cB_0/\delta)^{n/(c\eta)}\bigr)$, whereas the harmonic
condition forces $H+1>h_*\exp((B+1)/\delta)$.

(3) An exponent polynomial in budget and inverse tolerance, uniform in the cardinality of
the family and controlling the total positive excess, does not exist: by the exact cost
formula of \cref{sec:cost}, already in one variable the least admissible $H$ for families
of unrestricted cardinality satisfies $\log H=B/\delta+O_{p,h_0,\delta}(1)$ as
$B\to\infty$. The exponential dependence of \eqref{tors:eq:harmonicH} on budget over
tolerance is therefore unavoidable.

(4) Over one fixed finite extension $K/\Qp$ with ramification index $e$, a series in
$\OO_K[[T]]$ whose constant term is not a unit has valuation at least
$\min\{\rho,1/e\}$ at every witness of depth $\rho$, and a series with unit constant term
is nonzero on the whole open unit polydisc. A finite total budget then allows only
finitely many relevant series. The infinite-family statement in \cref{tors:thm:uniform}
has content because coefficients in $\OO_{\Cp}$, or in algebraic fields of unbounded
ramification, are allowed.
\end{remark}

